\section{Projective bundles. Ample sheaves}
\label{section:II.4}

\subsection{Definition of projective bundles.}
\label{subsection:II.4.1}

\begin{definition}[4.1.1]
\label{II.4.1.1}
Let $Y$ be a prescheme, $\mathcal{E}$ a quasi-coherent $\mathcal{O}_Y$-module,
and $\mathbf{S}_{\mathcal{O}_Y}(\mathcal{E})$ the symmetric
$\mathcal{O}_Y$-algebra of $\mathcal{E}$
(\hyperref[II.1.7.4]{1.7.4}), which is quasi-coherent
(\hyperref[II.1.7.7]{1.7.7}). We call the $Y$-scheme
$P=\operatorname{Proj}(\mathbf{S}_{\mathcal{O}_Y}(\mathcal{E}))$ the
\emph{projective bundle over $Y$ defined by $\mathcal{E}$}, and we denote it by
$\mathbf{P}(\mathcal{E})$. The
$\mathcal{O}_P$-module $\mathcal{O}_P(1)$ is called the \emph{tautological
sheaf on $P$}.
\end{definition}

When $Y$ is affine with ring $A$, we have $\mathcal{E}=\widetilde{E}$, where
$E$ is an $A$-module, and we then write $\mathbf{P}(E)$ instead of
$\mathbf{P}(\widetilde{E})$.

When we take $\mathcal{E}=\mathcal{O}_Y^n$, we write
$\mathbf{P}_Y^{n-1}$ instead of $\mathbf{P}(\mathcal{E})$; if moreover $Y$
is affine with ring $A$, we also write $\mathbf{P}_A^{n-1}$ instead of
$\mathbf{P}_Y^{n-1}$. Since
$\mathbf{S}_{\mathcal{O}_Y}(\mathcal{O}_Y)$ is canonically identified with
$\mathcal{O}_Y[T]$ (\hyperref[II.1.7.4]{1.7.4}),
$\mathbf{P}_Y^0$ is canonically identified with $Y$
(\hyperref[II.3.1.7]{3.1.7}); example
(\hyperref[II.2.4.3]{2.4.3}) is none other than $\mathbf{P}_K^1$.

\begin{env}[4.1.2]
\label{II.4.1.2}
Let $\mathcal{E}$, $\mathcal{F}$ be two quasi-coherent
$\mathcal{O}_Y$-modules, and $u:\mathcal{E}\to\mathcal{F}$ an
$\mathcal{O}_Y$-homomorphism; to it there corresponds canonically a
homomorphism
$\mathbf{S}(u):\mathbf{S}_{\mathcal{O}_Y}(\mathcal{E})\to
\mathbf{S}_{\mathcal{O}_Y}(\mathcal{F})$ of graded
$\mathcal{O}_Y$-algebras (\hyperref[II.1.7.4]{1.7.4}). If $u$ is
\emph{surjective}, so is $\mathbf{S}(u)$, and consequently
(\hyperref[II.3.6.2]{3.6.2, (i)}),
$\operatorname{Proj}(\mathbf{S}(u))$ is a closed immersion
$\mathbf{P}(\mathcal{F})\to\mathbf{P}(\mathcal{E})$, which we will denote by
$\mathbf{P}(u)$. We can therefore say that $\mathbf{P}(\mathcal{E})$ is a
\emph{contravariant functor in $\mathcal{E}$}, provided that we take as
morphisms of quasi-coherent $\mathcal{O}_Y$-modules only the surjective
homomorphisms.

Still supposing $u$ surjective, and setting
$P=\mathbf{P}(\mathcal{E})$, $Q=\mathbf{P}(\mathcal{F})$ and
$j=\mathbf{P}(u)$, we have, up to isomorphism,
\[
\label{II.4.1.2.1}
  j^*(\mathcal{O}_P(n))=\mathcal{O}_Q(n)
  \qquad\text{for every }n\in\mathbf{Z}
\tag{4.1.2.1}
\]
as follows from (\hyperref[II.3.6.3]{3.6.3}).
\end{env}

\begin{env}[4.1.3]
\label{II.4.1.3}
Now let $\psi:Y'\to Y$ be a morphism, and let
$\mathcal{E}'=\psi^*(\mathcal{E})$; we then have
$\mathbf{S}_{\mathcal{O}_{Y'}}(\mathcal{E}')=
\psi^*(\mathbf{S}_{\mathcal{O}_Y}(\mathcal{E}))$
(\hyperref[II.1.7.5]{1.7.5}); consequently
(\hyperref[II.3.5.3]{3.5.3}), we have
\[
\label{II.4.1.3.1}
  \mathbf{P}(\psi^*(\mathcal{E}))=
  \mathbf{P}(\mathcal{E})\times_YY'
\tag{4.1.3.1}
\]
up to canonical isomorphism; moreover, we evidently have
\[
  \psi^*((\mathbf{S}_{\mathcal{O}_Y}(\mathcal{E}))(n))
  =(\mathbf{S}_{\mathcal{O}_{Y'}}(\mathcal{E}'))(n)
\]
for every $n\in\mathbf{Z}$, hence, setting $P=\mathbf{P}(\mathcal{E})$,
$P'=\mathbf{P}(\mathcal{E}')$, we have
(\hyperref[II.3.5.4]{3.5.4}), up to isomorphism,
\[
\label{II.4.1.3.2}
  \mathcal{O}_{P'}(n)=\mathcal{O}_P(n)\otimes_Y\mathcal{O}_{Y'}
  \qquad\text{for every }n\in\mathbf{Z}.
\tag{4.1.3.2}
\]
\end{env}

\oldpage[II]{72}

\begin{proposition}[4.1.4]
\label{II.4.1.4}
Let $\mathcal{L}$ be an invertible $\mathcal{O}_Y$-module. For every
quasi-coherent $\mathcal{O}_Y$-module $\mathcal{E}$, there exists a canonical
$Y$-isomorphism
$i_{\mathcal{L}}:\mathbf{P}(\mathcal{E})\xrightarrow{\sim}
\mathbf{P}(\mathcal{E}\otimes\mathcal{L})$; moreover, if we set
$P=\mathbf{P}(\mathcal{E})$, $Q=\mathbf{P}(\mathcal{E}\otimes\mathcal{L})$,
then $i_{\mathcal{L}}^*(\mathcal{O}_Q(n))$ is canonically isomorphic to
$\mathcal{O}_P(n)\otimes_Y\mathcal{L}^{\otimes n}$ for every
$n\in\mathbf{Z}$.
\end{proposition}

\begin{proof}
Let us first remark that if $A$ is a ring, $E$ an $A$-module, and $L$ a
\emph{free monogenic} $A$-module, we canonically define a homomorphism
of $A$-modules
\[
  \mathbf{S}_n(E\otimes L)\longrightarrow
  \mathbf{S}_n(E)\otimes L^{\otimes n}
\]
by assigning to
$(x_1\otimes y_1)\cdots(x_n\otimes y_n)$ the element
\[
  (x_1x_2\cdots x_n)\otimes(y_1\otimes y_2\otimes\cdots\otimes y_n)
  \qquad(x_i\in E,\ y_i\in L\text{ for }1\leq i\leq n)~;
\]
we verify immediately (by reducing to the case where $L=A$) that this
homomorphism is in fact an isomorphism. From this we deduce a canonical
isomorphism of graded $A$-algebras
$\mathbf{S}_A(E\otimes L)\xrightarrow{\sim}
\bigoplus_{n\geq0}\mathbf{S}_n(E)\otimes L^{\otimes n}$. Returning to the
conditions of (4.1.4), the preceding remarks allow us to define a
canonical isomorphism of graded $\mathcal{O}_Y$-algebras
\[
\label{II.4.1.4.1}
  \mathbf{S}_{\mathcal{O}_Y}(\mathcal{E}\otimes_{\mathcal{O}_Y}\mathcal{L})
  \xrightarrow{\sim}
  \bigoplus_{n\geq0}\mathbf{S}_n(\mathcal{E})
  \otimes_{\mathcal{O}_Y}\mathcal{L}^{\otimes n}
\tag{4.1.4.1}
\]
by defining this isomorphism as an isomorphism of presheaves and
taking into account (\hyperref[II.1.7.4]{1.7.4}),
(\hyperref[I.1.3.9]{I, 1.3.9}) and
(\hyperref[I.1.3.12]{I, 1.3.12}). The proposition then follows from
(\hyperref[II.3.1.8]{3.1.8, (iii)}) and
(\hyperref[II.3.2.10]{3.2.10}).
\end{proof}

\begin{env}[4.1.5]
\label{II.4.1.5}
With the hypotheses of (\hyperref[II.4.1.1]{4.1.1}), set
$P=\mathbf{P}(\mathcal{E})$, and denote by $p$ the structural morphism
$P\to Y$. Since by definition we have
$\mathcal{E}=(\mathbf{S}_{\mathcal{O}_Y}(\mathcal{E}))_1$, we have a
canonical homomorphism
$\alpha_1:\mathcal{E}\to p_*(\mathcal{O}_P(1))$
(\hyperref[II.3.3.2.2]{3.3.2.2}), and consequently also
(\hyperref[0.4.4.3]{0, 4.4.3}) a canonical homomorphism
\[
\label{II.4.1.5.1}
  \alpha_1^\sharp:p^*(\mathcal{E})\longrightarrow\mathcal{O}_P(1).
\tag{4.1.5.1}
\]
\end{env}

\begin{proposition}[4.1.6]
\label{II.4.1.6}
The canonical homomorphism (\hyperref[II.4.1.5.1]{4.1.5.1}) is surjective.
\end{proposition}

\begin{proof}
Indeed, we saw in (\hyperref[II.3.3.2]{3.3.2}) that
$\alpha_1^\sharp$ corresponds functorially to the canonical homomorphism
\[
  \mathcal{E}\otimes_{\mathcal{O}_Y}
  \mathbf{S}_{\mathcal{O}_Y}(\mathcal{E})
  \longrightarrow
  (\mathbf{S}_{\mathcal{O}_Y}(\mathcal{E}))(1)~;
\]
since by definition $\mathcal{E}$ generates
$\mathbf{S}_{\mathcal{O}_Y}(\mathcal{E})$, this homomorphism is surjective,
whence the conclusion by virtue of (\hyperref[II.3.2.4]{3.2.4}).
\end{proof}

\subsection{Morphisms from a prescheme to a projective bundle.}
\label{subsection:II.4.2}

\begin{env}[4.2.1]
\label{II.4.2.1}
Keeping the notation of (\hyperref[II.4.1.5]{4.1.5}), let $X$ be a
$Y$-prescheme, $q:X\to Y$ the structural morphism, and let $r:X\to P$ be a
$Y$-\emph{morphism}, so that we have the commutative diagram
\[
  \xymatrix{
    P \ar[d]_p & X \ar[l]_r \ar[dl]^q
    \\Y
  }
\]

\oldpage[II]{73}

Since the functor $r^*$ is right exact
(\hyperref[0.4.3.1]{0, 4.3.1}), we deduce from the homomorphism
(\hyperref[II.4.1.5.1]{4.1.5.1}), which is surjective
(\hyperref[II.4.1.6]{4.1.6}), a surjective homomorphism
\[
  r^*(\alpha_1^\sharp):r^*(p^*(\mathcal{E}))\longrightarrow
  r^*(\mathcal{O}_P(1)).
\]
But $r^*(p^*(\mathcal{E}))=q^*(\mathcal{E})$, and
$r^*(\mathcal{O}_P(1))$ is locally isomorphic to
$r^*(\mathcal{O}_P)=\mathcal{O}_X$; in other words, it is an invertible
sheaf $\mathcal{L}_r$ on $\mathcal{O}_X$, and we have thus defined, starting
from $r$, a canonical surjective $\mathcal{O}_X$-homomorphism
\[
\label{II.4.2.1.1}
  \varphi_r:q^*(\mathcal{E})\longrightarrow\mathcal{L}_r.
\tag{4.2.1.1}
\]

When $Y=\operatorname{Spec}(A)$ is affine and
$\mathcal{E}=\widetilde{E}$, we can again make this homomorphism explicit as
follows: given $f\in E$, it first follows from
(\hyperref[II.2.6.3]{2.6.3}) that we have
\[
\label{II.4.2.1.2}
  r^{-1}(D_+(f))=X_{\varphi_r^\flat(f)}.
\tag{4.2.1.2}
\]
On the other hand, let $V$ be an affine open subset of $X$ contained in
$r^{-1}(D_+(f))$, and let $B$ be its ring, which is an $A$-algebra; set
$S=\mathbf{S}_A(E)$; the restriction of $r$ to $V$ corresponds to an
$A$-homomorphism $\omega:S_{(f)}\to B$, and we have
$q^*(\mathcal{E})|V=\widetilde{E\otimes_A B}$ and
$\mathcal{L}_r|V=\widetilde{L_r}$, where
$L_r=(S(1))_{(f)}\otimes_{S_{(f)}}B_{[\omega]}$
(\hyperref[I.1.6.5]{I, 1.6.5}). The restriction of $\varphi_r$ to
$q^*(\mathcal{E})|V$ corresponds to the $B$-homomorphism
$u:E\otimes_A B\to L_r$ that sends $x\otimes1$ to
$(x/1)\otimes1=(f/1)\otimes\omega(x/f)$. The canonical extension of
$\varphi_r$ to a homomorphism of $\mathcal{O}_X$-algebras
\[
  \psi_r:q^*(\mathbf{S}(\mathcal{E}))=
  \mathbf{S}(q^*(\mathcal{E}))\longrightarrow
  \mathbf{S}(\mathcal{L}_r)=\bigoplus_{n\geq0}\mathcal{L}_r^{\otimes n}
\]
is therefore such that the restriction of $\psi_r$ to
$q^*(\mathbf{S}_n(\mathcal{E}))|V$ corresponds to the homomorphism
$\mathbf{S}_n(E\otimes_A B)=\mathbf{S}_n(E)\otimes_A B\to L_r^{\otimes n}$
that sends $s\otimes1$ to
$(f/1)^{\otimes n}\otimes\omega(s/f^n)$.
\end{env}

\begin{env}[4.2.2]
\label{II.4.2.2}
Conversely, suppose given a morphism $q:X\to Y$, an invertible
$\mathcal{O}_X$-module $\mathcal{L}$ and a quasi-coherent
$\mathcal{O}_Y$-module $\mathcal{E}$; to a homomorphism
$\varphi:q^*(\mathcal{E})\to\mathcal{L}$ there corresponds canonically a
homomorphism of quasi-coherent $\mathcal{O}_X$-algebras
\[
  \psi:\mathbf{S}(q^*(\mathcal{E}))=q^*(\mathbf{S}(\mathcal{E}))
  \longrightarrow\bigoplus_{n\geq0}\mathcal{L}^{\otimes n},
\]
and consequently (\hyperref[II.3.7.1]{3.7.1}) a $Y$-morphism
$r_{\mathcal{L},\psi}:G(\psi)\to
\operatorname{Proj}(\mathbf{S}(\mathcal{E}))=\mathbf{P}(\mathcal{E})$, which
we will denote by $r_{\mathcal{L},\varphi}$. If $\varphi$ is surjective, so
is $\psi$, hence (\hyperref[II.3.7.4]{3.7.4})
$r_{\mathcal{L},\varphi}$ is everywhere defined. Moreover, with the notation of
(\hyperref[II.4.2.1]{4.2.1}) and of (4.2.2):
\end{env}

\begin{proposition}[4.2.3]
\label{II.4.2.3}
Given a morphism $q:X\to Y$ and a quasi-coherent $\mathcal{O}_Y$-module
$\mathcal{E}$, the maps
$r\mapsto(\mathcal{L}_r,\varphi_r)$ and
$(\mathcal{L},\varphi)\mapsto r_{\mathcal{L},\varphi}$ put in
bijective correspondence the set of $Y$-morphisms
$r:X\to\mathbf{P}(\mathcal{E})$ and the set of equivalence classes of
pairs $(\mathcal{L},\varphi)$ consisting of an invertible $\mathcal{O}_X$-module
$\mathcal{L}$ and a surjective homomorphism
$\varphi:q^*(\mathcal{E})\to\mathcal{L}$, two pairs
$(\mathcal{L},\varphi)$, $(\mathcal{L}',\varphi')$ being equivalent if there
exists an $\mathcal{O}_X$-isomorphism
$\tau:\mathcal{L}\xrightarrow{\sim}\mathcal{L}'$ such that
$\varphi'=\tau\circ\varphi$.
\end{proposition}

\begin{proof}
Let us first start from a $Y$-morphism $r:X\to\mathbf{P}(\mathcal{E})$, form
$\mathcal{L}_r$ and $\varphi_r$ (\hyperref[II.4.2.1]{4.2.1}), then
$r'=r_{\mathcal{L}_r,\varphi_r}$; it follows immediately from
(\hyperref[II.4.2.1]{4.2.1}) and from
(\hyperref[II.3.7.2]{3.7.2}) that the morphisms $r$ and $r'$ are
identical (by taking as generator of $\mathcal{L}_r$ the element
$(f/1)\otimes1$ in order to define the homomorphisms $v_n$ of
(\hyperref[II.3.7.2]{3.7.2})). Conversely, let us start from a pair
$(\mathcal{L},\varphi)$ and form

\oldpage[II]{74}

$r''=r_{\mathcal{L},\varphi}$, then $\mathcal{L}_{r''}$ and
$\varphi_{r''}$; let us show that there is a canonical isomorphism
$\tau:\mathcal{L}_{r''}\xrightarrow{\sim}\mathcal{L}$ such that
$\varphi=\tau\circ\varphi_{r''}$; to define it we may place ourselves in the
case where $Y=\operatorname{Spec}(A)$ and $X=\operatorname{Spec}(B)$ are affine,
and (with the notation of (\hyperref[II.4.2.1]{4.2.1}) and of
(\hyperref[II.3.7.2]{3.7.2})) assign to every element
$(x/1)\otimes1$ of $L_{r''}$ (with $x\in E$) the element $v_1(x)c$ of $L$. We
verify at once that $\tau$ does not depend on the chosen generator $c$ of $L$;
since $v_1$ is surjective by hypothesis, to prove that $\tau$ is an
isomorphism it suffices to see that if $x/1=0$ in $(S(1))_{(f)}$, we have
$v_1(x)/1=0$ in $B_g$; but the first relation means that $f^nx=0$
in $\mathbf{S}_{n+1}(E)$ for some $n$, and we deduce from this that
$v_{n+1}(f^nx)=g^nv_1(x)=0$ in $B$, whence the conclusion. Finally, it is
immediate that for two equivalent pairs $(\mathcal{L},\varphi)$,
$(\mathcal{L}',\varphi')$, we have
$r_{\mathcal{L},\varphi}=r_{\mathcal{L}',\varphi'}$.
\end{proof}

In particular, for $X=Y$:

\begin{theorem}[4.2.4]
\label{II.4.2.4}
The set of $Y$-sections of $\mathbf{P}(\mathcal{E})$ is in canonical
bijective correspondence with the set of quasi-coherent
sub-$\mathcal{O}_Y$-modules $\mathcal{F}$ of $\mathcal{E}$
such that $\mathcal{E}/\mathcal{F}$ is invertible.
\end{theorem}

We observe that this property agrees with the classical definition of
``projective space'' as the set of hyperplanes of a vector
space (the classical situation being the case $Y=\operatorname{Spec}(K)$, where
$K$ is a field, and $\mathcal{E}=\widetilde{E}$, $E$ being a finite-dimensional
$K$-vector space; the $\mathcal{F}$ having the property stated
in (\hyperref[II.4.2.4]{4.2.4}) then correspond to the hyperplanes of
$E$, and we know on the other hand that the $Y$-sections of
$\mathbf{P}(\mathcal{E})$ are then the points of
$\mathbf{P}(\mathcal{E})$ rational over $K$
(\hyperref[I.3.4.5]{I, 3.4.5})).

\begin{remark}[4.2.5]
\label{II.4.2.5}
Since there is a canonical bijective correspondence between $Y$-morphisms from $X$
to $P$ and their graph morphisms, which are $X$-sections of $P\times_YX$
(\hyperref[I.3.3.14]{I, 3.3.14}), we see that, conversely,
(\hyperref[II.4.2.3]{4.2.3}) can be deduced from
(\hyperref[II.4.2.4]{4.2.4}). Let us denote by
$\operatorname{Hyp}_Y(X,\mathcal{E})$ the set of quasi-coherent
sub-$\mathcal{O}_X$-modules $\mathcal{F}$ of
$\mathcal{E}\otimes_Y\mathcal{O}_X=q^*(\mathcal{E})$ such that
$q^*(\mathcal{E})/\mathcal{F}$ is an invertible $\mathcal{O}_X$-module. If
$g:X'\to X$ is a $Y$-morphism, it follows from the fact that $g^*$ is right
exact that we have
\[
  g^*(q^*(\mathcal{E})/\mathcal{F})=
  g^*(q^*(\mathcal{E}))/g^*(\mathcal{F}),
\]
hence that the latter sheaf is invertible, and consequently
$\operatorname{Hyp}_Y(X,\mathcal{E})$ is a contravariant functor on the
category of $Y$-preschemes. We can then interpret Theorem
(\hyperref[II.4.2.4]{4.2.4}) as defining a canonical isomorphism
between the functors $\operatorname{Hom}_Y(X,\mathbf{P}(\mathcal{E}))$ and
$\operatorname{Hyp}_Y(X,\mathcal{E})$, contravariant in $X$ varying in the
category of $Y$-preschemes. This also provides a characterization of the
projective bundle $P=\mathbf{P}(\mathcal{E})$ by the following universal property,
closer to geometric intuition than the constructions of \S\S~2 and
3: for every morphism $q:X\to Y$ and every invertible $\mathcal{O}_X$-module
$\mathcal{L}$ that is a quotient of $\mathcal{E}\otimes_{\mathcal{O}_Y}\mathcal{O}_X$,
there exists a unique $Y$-morphism $r:X\to P$ such that
$\mathcal{L}=r^*(\mathcal{O}_P(1))$.

We will see later how one can define in the same way, among others,
the ``Grassmannian'' schemes.
\end{remark}

\begin{corollary}[4.2.6]
\label{II.4.2.6}
Suppose that every invertible $\mathcal{O}_Y$-module is trivial
(\hyperref[I.2.4.8]{I, 2.4.8}). Let $V$ be the group
$\operatorname{Hom}_{\mathcal{O}_Y}(\mathcal{E},\mathcal{O}_Y)$, considered
as a module over the ring $A=\Gamma(Y,\mathcal{O}_Y)$, and let $V^*$ be the
subset of $V$ consisting of the surjective homomorphisms. Then the set
of $Y$-sections of $\mathbf{P}(\mathcal{E})$ is canonically identified with
$V^*/A^*$, where $A^*$ is the group of units of $A$.
\end{corollary}

\oldpage[II]{75}

In particular:
\begin{enumerate}
  \item[1\textsuperscript{o}] Corollary
  (\hyperref[II.4.2.6]{4.2.6}) applies when $Y$ is a
  \emph{local scheme} (\hyperref[I.2.4.8]{I, 2.4.8}). Let $Y$ be an
  arbitrary prescheme, $y$ a point of $Y$, and
  $Y'=\operatorname{Spec}(k(y))$; the fibre $p^{-1}(y)$ of
  $\mathbf{P}(\mathcal{E})$ is, by virtue of
  (\hyperref[II.4.1.3.1]{4.1.3.1}), identified with
  $\mathbf{P}(\mathcal{E}^y)$, where
  $\mathcal{E}^y=\mathcal{E}_y\otimes_{\mathcal{O}_y}k(y)
  =\mathcal{E}_y/\mathfrak{m}_y\mathcal{E}_y$ is considered as a vector
  space over $k(y)$. More generally, if $K$ is an extension of
  $k(y)$, $p^{-1}(y)\otimes_{k(y)}K$ is identified with
  $\mathbf{P}(\mathcal{E}^y\otimes_{k(y)}K)$. Corollary
  (\hyperref[II.4.2.6]{4.2.6}) therefore shows that the set of
  \emph{geometric points of $\mathbf{P}(\mathcal{E})$ with values in
  the extension $K$ of $k(y)$} (\hyperref[I.3.4.5]{I, 3.4.5}), which one
  can also call the \emph{rational geometric fibre over $K$ of
  $\mathbf{P}(\mathcal{E})$ above the point $y$}, is identified with
  the \emph{projective space} associated to the \emph{dual} of the $K$-vector
  space $\mathcal{E}^y\otimes_{k(y)}K$.
  \item[2\textsuperscript{o}] Suppose that $Y$ is affine with ring $A$, and
  moreover that every invertible $\mathcal{O}_Y$-module is trivial; take,
  moreover, $\mathcal{E}=\mathcal{O}_Y^n$; then, in
  (\hyperref[II.4.2.6]{4.2.6}), $V$ is identified with $A^n$
  (\hyperref[I.1.3.8]{I, 1.3.8}), and $V^*$ with the set of systems
  $(f_i)_{1\leq i\leq n}$ of elements of $A$ generating the ideal $A$; two
  such systems define the same $Y$-section of
  $\mathbf{P}_Y^{n-1}=\mathbf{P}_A^{n-1}$, in other words the same point of
  $\mathbf{P}_A^{n-1}$ with values in $A$, if and only if one is deduced
  from the other by multiplication by an invertible element of $A$.
\end{enumerate}

These properties justify the terminology ``projective bundle'' for
$\mathbf{P}(\mathcal{E})$. Note that the definition of ``projective
space'' that we thus obtain is in fact dual to the classical
definition; this is imposed on us by the necessity of being able to define
$\mathbf{P}(\mathcal{E})$ for an arbitrary quasi-coherent $\mathcal{O}_Y$-module
$\mathcal{E}$, not necessarily locally free.

\begin{remark}[4.2.7]
\label{II.4.2.7}
We will see in Chapter~V that if $Y$ is locally Noetherian, connected and non-empty,
and if $\mathcal{E}$ is locally free of rank at least $2$, then, setting
$P=\mathbf{P}(\mathcal{E})$, every invertible $\mathcal{O}_P$-module
$\mathcal{L}$ is isomorphic to an $\mathcal{O}_P$-module of the form
$\mathcal{L}'\otimes_{\mathcal{O}_Y}\mathcal{O}_P(m)$, where $\mathcal{L}'$ is
an invertible $\mathcal{O}_Y$-module, uniquely determined up to isomorphism,
and $m$ a uniquely determined integer. In other words,
$\mathrm{H}^1(P,\mathcal{O}_P^*)$ is isomorphic to
$\mathbf{Z}\times\mathrm{H}^1(Y,\mathcal{O}_Y^*)$
(\hyperref[0.5.4.7]{0, 5.4.7}). We will also see
(\hyperref[III.2.1.15]{III, 2.1.15}, taking into account
(\hyperref[0.5.4.10]{0, 5.4.10})) that
$p_*(\mathcal{L})=0$ if $m<0$ and that
$p_*(\mathcal{L})$ is isomorphic to
$\mathcal{L}'\otimes_{\mathcal{O}_Y}
(\mathbf{S}_{\mathcal{O}_Y}(\mathcal{E}))_m$ if $m\geq0$. If $\mathcal{F}$
is a quasi-coherent $\mathcal{O}_Y$-module, every $Y$-morphism
$\mathbf{P}(\mathcal{E})\to\mathbf{P}(\mathcal{F})$ is therefore determined by
the data of an invertible $\mathcal{O}_Y$-module $\mathcal{L}'$, an integer
$m\geq0$ and an $\mathcal{O}_Y$-homomorphism
$\psi:\mathcal{F}\to\mathcal{L}'\otimes_{\mathcal{O}_Y}
(\mathbf{S}_{\mathcal{O}_Y}(\mathcal{E}))_m$ such that the
corresponding homomorphism $\psi^\sharp$ of $\mathcal{O}_{\mathbf{P}(\mathcal{E})}$-modules
is surjective. We will also see that if the $Y$-morphism considered is an
isomorphism, then $m=1$ and $\mathcal{F}$ is isomorphic to
$\mathcal{E}\otimes_{\mathcal{O}_Y}\mathcal{L}'$ (converse of
(\hyperref[II.4.1.4]{4.1.4})). This will make it possible to determine the
sheaf of germs of automorphisms of $\mathbf{P}(\mathcal{E})$ as the quotient
of the sheaf of groups $\mathcal{A}ut(\mathcal{E})$ (locally isomorphic to
$\operatorname{GL}(n,\mathcal{O}_Y)$ if $\mathcal{E}$ is of rank $n$) by
$\mathcal{O}_Y^*$.
\end{remark}

\begin{env}[4.2.8]
\label{II.4.2.8}
Keeping the notation of (\hyperref[II.4.2.1]{4.2.1}), let
$u:X'\to X$ be a morphism; if the $Y$-morphism $r:X\to P$ corresponds to
the homomorphism $\varphi:q^*(\mathcal{E})\to\mathcal{L}$, the $Y$-morphism
$r\circ u$ corresponds to
$u^*(\varphi):u^*(q^*(\mathcal{E}))\to u^*(\mathcal{L})$, as follows
immediately from the definitions.
\end{env}

\begin{env}[4.2.9]
\label{II.4.2.9}
Let $\mathcal{E}$, $\mathcal{F}$ be two quasi-coherent $\mathcal{O}_Y$-modules,
$v:\mathcal{E}\to\mathcal{F}$ a surjective homomorphism, and
$j=\mathbf{P}(v)$ the corresponding closed immersion
$\mathbf{P}(\mathcal{F})\to\mathbf{P}(\mathcal{E})$
(\hyperref[II.4.1.2]{4.1.2}). If the $Y$-morphism
$r:X\to\mathbf{P}(\mathcal{F})$ corresponds to the homomorphism
$\varphi:q^*(\mathcal{F})\to\mathcal{L}$, the

\oldpage[II]{76}

$Y$-morphism $j\circ r$ corresponds to
\[
  q^*(\mathcal{E})\xrightarrow{q^*(v)}q^*(\mathcal{F})
  \xrightarrow{\varphi}\mathcal{L}~;
\]
this again follows from the definition given in
(\hyperref[II.4.2.1]{4.2.1}).
\end{env}

\begin{env}[4.2.10]
\label{II.4.2.10}
Let $\psi:Y'\to Y$ be a morphism, and set
$\mathcal{E}'=\psi^*(\mathcal{E})$. If the $Y$-morphism $r:X\to P$ corresponds
to the homomorphism $\varphi:q^*(\mathcal{E})\to\mathcal{L}$, the
$Y'$-morphism
\[
  r_{(Y')}:X_{(Y')}\longrightarrow P'=\mathbf{P}(\mathcal{E}')
\]
corresponds to
$\varphi_{(Y')}:q_{(Y')}^*(\mathcal{E}')=
  q^*(\mathcal{E})\otimes_Y\mathcal{O}_{Y'}\to
\mathcal{L}\otimes_Y\mathcal{O}_{Y'}$. Indeed, by virtue of
(\hyperref[II.4.1.3.1]{4.1.3.1}), we have the commutative diagram
\[
  \xymatrix{
    Y' \ar[d]
    & P'=P_{(Y')} \ar[l]_{p_{(Y')}} \ar[d]^{u}
    & X_{(Y')} \ar[l]_{r_{(Y')}} \ar[d]^{v}
    \\Y
    & P \ar[l]_{p}
    & X \ar[l]_{r}
  }
\]
By (\hyperref[II.4.1.3.2]{4.1.3.2}), we have
\[
  (r_{(Y')})^*(\mathcal{O}_{P'}(1))
  =(r_{(Y')})^*(u^*(\mathcal{O}_P(1)))
  =v^*(r^*(\mathcal{O}_P(1)))
  =v^*(\mathcal{L})=\mathcal{L}\otimes_Y\mathcal{O}_{Y'}~;
\]
on the other hand, $u^*(\alpha_1^\sharp)$ is none other than the canonical
homomorphism
$\alpha_1^\sharp:(p_{(Y')})^*(\mathcal{E}')\to\mathcal{O}_{P'}(1)$; this
is seen by making the canonical homomorphisms $\alpha_1^\sharp$ on
$P$ and $P'$ explicit as in (\hyperref[II.4.1.6]{4.1.6}). Whence our assertion.
\end{env}

\subsection{The Segre morphism.}
\label{subsection:II.4.3}

\begin{env}[4.3.1]
\label{II.4.3.1}
Let $Y$ be a prescheme, and $\mathcal{E}$, $\mathcal{F}$ two quasi-coherent
$\mathcal{O}_Y$-modules; set
$P_1=\mathbf{P}(\mathcal{E})$, $P_2=\mathbf{P}(\mathcal{F})$, and denote
by $p_1:P_1\to Y$, $p_2:P_2\to Y$ the structural morphisms. Let
$Q=P_1\times_YP_2$, and let $q_1:Q\to P_1$, $q_2:Q\to P_2$ be the
canonical projections; the $\mathcal{O}_Q$-module
\[
  \mathcal{L}=\mathcal{O}_{P_1}(1)\otimes_Y\mathcal{O}_{P_2}(1)
  :=q_1^*(\mathcal{O}_{P_1}(1))\otimes_{\mathcal{O}_Q}
  q_2^*(\mathcal{O}_{P_2}(1))
\]
is invertible, as the tensor product of two invertible
$\mathcal{O}_Q$-modules (\hyperref[0.5.4.4]{0, 5.4.4}). On the other hand, if
$r=p_1\circ q_1=p_2\circ q_2$ is the structural morphism $Q\to Y$, we have
\[
  r^*(\mathcal{E}\otimes_{\mathcal{O}_Y}\mathcal{F})
  =q_1^*(p_1^*(\mathcal{E}))\otimes_{\mathcal{O}_Q}
  q_2^*(p_2^*(\mathcal{F}))
\]
(\hyperref[0.4.3.3]{0, 4.3.3}); the canonical surjective homomorphisms
(\hyperref[II.4.1.5.1]{4.1.5.1})
$p_1^*(\mathcal{E})\to\mathcal{O}_{P_1}(1)$,
$p_2^*(\mathcal{F})\to\mathcal{O}_{P_2}(1)$ therefore give, by tensor
product, a canonical homomorphism
\[
\label{II.4.3.1.1}
  s:r^*(\mathcal{E}\otimes_{\mathcal{O}_Y}\mathcal{F})\longrightarrow
  \mathcal{L}
\tag{4.3.1.1}
\]
which is evidently surjective; we therefore deduce from it
(\hyperref[II.4.2.2]{4.2.2}) a canonical morphism, called the
\emph{Segre morphism}:
\[
\label{II.4.3.1.2}
  \varsigma:\mathbf{P}(\mathcal{E})\times_Y\mathbf{P}(\mathcal{F})
  \longrightarrow
  \mathbf{P}(\mathcal{E}\otimes_{\mathcal{O}_Y}\mathcal{F}).
\tag{4.3.1.2}
\]
Let us make the morphism $\varsigma$ explicit when $Y=\operatorname{Spec}(A)$ is
affine, $\mathcal{E}=\widetilde{E}$, $\mathcal{F}=\widetilde{F}$, $E$ and $F$
being two $A$-modules, whence
$\mathcal{E}\otimes_{\mathcal{O}_Y}\mathcal{F}=\widetilde{E\otimes_A F}$
(\hyperref[I.1.3.12]{I, 1.3.12}); set
$R=\mathbf{S}_A(E)$, $S=\mathbf{S}_A(F)$,
$T=\mathbf{S}_A(E\otimes_A F)$; let $f\in E$, $g\in F$, and consider
the affine open subset
\[
  D_+(f)\times_YD_+(g)=\operatorname{Spec}(B)
\]

\oldpage[II]{77}

of $Q$, where $B=R_{(f)}\otimes_AS_{(g)}$; the restriction of $\mathcal{L}$ to
this affine open subset is $\widetilde{L}$, where
\[
  L=(R(1))_{(f)}\otimes_A(S(1))_{(g)}
\]
and the element $c=(f/1)\otimes(g/1)$ is a generator of $L$ considered as a
free $B$-module (\hyperref[II.2.5.7]{2.5.7}). The homomorphism
(\hyperref[II.4.3.1.1]{4.3.1.1}) corresponds to the homomorphism
\[
  (x\otimes y)\otimes b\longmapsto b((x/1)\otimes(y/1))
\]
from $(E\otimes_A F)\otimes_AB$ to $L$. With the notation of
(\hyperref[II.3.7.2]{3.7.2}), we therefore have
$v_1(x\otimes y)=(x/f)\otimes(y/g)$; the restriction of $\varsigma$ to
$D_+(f)\times_YD_+(g)$ is a morphism from this affine scheme to
$D_+(f\otimes g)$, corresponding to the ring homomorphism
$\omega:T_{(f\otimes g)}\to R_{(f)}\otimes_AS_{(g)}$ such that
\[
\label{II.4.3.1.3}
  \omega((x\otimes y)/(f\otimes g))=(x/f)\otimes(y/g)
\tag{4.3.1.3}
\]
for $x\in E$ and $y\in F$.
\end{env}

\begin{env}[4.3.2]
\label{II.4.3.2}
It follows from (\hyperref[II.4.2.3]{4.2.3}) that we have a canonical
isomorphism
\[
\label{II.4.3.2.1}
  \tau:\varsigma^*(\mathcal{O}_P(1))\xrightarrow{\sim}
  \mathcal{O}_{P_1}(1)\otimes_Y\mathcal{O}_{P_2}(1)
\tag{4.3.2.1}
\]
where we have set
$P=\mathbf{P}(\mathcal{E}\otimes_{\mathcal{O}_Y}\mathcal{F})$. Let us show
that, for $x\in\Gamma(Y,\mathcal{E})$ and
$y\in\Gamma(Y,\mathcal{F})$, we have
\[
\label{II.4.3.2.2}
  \tau(\alpha_1(x\otimes y))=\alpha_1(x)\otimes\alpha_1(y).
\tag{4.3.2.2}
\]
Indeed, we are reduced to the case where $Y$ is affine, and we then have, with
the notation of (\hyperref[II.4.3.1]{4.3.1}) and of
(\hyperref[II.2.6.2]{2.6.2}),
$\alpha_1^{f\otimes g}(x\otimes y)=(x\otimes y)/1$ in
$(T(1))_{(f\otimes g)}$, $\alpha_1^f(x)=x/1$ in $(R(1))_{(f)}$ and
$\alpha_1^g(y)=y/1$ in $(S(1))_{(g)}$. The definition of $\tau$ given in
(\hyperref[II.4.2.3]{4.2.3}) and the computation of $v_1$ made in
(\hyperref[II.4.3.1]{4.3.1}) then immediately prove
(\hyperref[II.4.3.2.2]{4.3.2.2}). From this we deduce the formula
\[
\label{II.4.3.2.3}
  \varsigma^{-1}(P_{x\otimes y})=(P_1)_x\times_Y(P_2)_y
\tag{4.3.2.3}
\]
with the notation of (\hyperref[II.3.1.4]{3.1.4}). Indeed, taking into account
(\hyperref[II.3.3.3]{3.3.3}), formula
(\hyperref[II.4.3.2.2]{4.3.2.2}) reduces us (by returning to the affine case,
with the help of (\hyperref[I.3.2.7]{I, 3.2.7} and
(\hyperref[I.3.2.3]{3.2.3})) to proving the following lemma:

\begin{lemma}[4.3.2.4]
\label{II.4.3.2.4}
Let $B$, $B'$ be two $A$-algebras, and let
$Y=\operatorname{Spec}(A)$, $Z=\operatorname{Spec}(B)$,
$Z'=\operatorname{Spec}(B')$; for all $t\in B$, $t'\in B'$, we have
$D(t\otimes t')=D(t)\times_YD(t')$.
\end{lemma}

\begin{proof}
Indeed, if $p:Z\times_YZ'\to Z$ and $p':Z\times_YZ'\to Z'$ are the
canonical projections, it follows from
(\hyperref[I.1.2.2.2]{I, 1.2.2.2}) that we have
$p^{-1}(D(t))=D(t\otimes1)$ and
$p'^{-1}(D(t'))=D(1\otimes t')$; the conclusion follows from
(\hyperref[I.3.2.7]{I, 3.2.7}) and
(\hyperref[I.1.1.9.1]{I, 1.1.9.1}), since
$(t\otimes1)(1\otimes t')=t\otimes t'$.
\end{proof}
\end{env}

\begin{proposition}[4.3.3]
\label{II.4.3.3}
The Segre morphism is a closed immersion.
\end{proposition}

\begin{proof}
Indeed, the question being local on $Y$, we are reduced to the case where $Y$ is
affine. With the notation of (\hyperref[II.4.3.1]{4.3.1}) and
(\hyperref[II.4.3.2]{4.3.2}), the $D_+(f\otimes g)$ then form a basis
of the topology of $P$, since the $f\otimes g$ generate $T$ when $f$
ranges over $E$ and $g$ ranges over $F$. On the other hand, by virtue of
(\hyperref[II.4.3.2.3]{4.3.2.3}) we have
$\varsigma^{-1}(D_+(f\otimes g))=D_+(f)\times_YD_+(g)$. It therefore suffices
(\hyperref[I.4.2.4]{I, 4.2.4}) to prove that the restriction of
$\varsigma$ to $D_+(f)\times_YD_+(g)$ is a closed immersion into
$D_+(f\otimes g)$. Now, with the same notation, formula
(\hyperref[II.4.3.1.3]{4.3.1.3}) shows that $\omega$ is surjective, which
completes the proof.
\end{proof}

\begin{env}[4.3.4]
\label{II.4.3.4}
The Segre morphism is functorial in $\mathcal{E}$ and $\mathcal{F}$,
when one restricts the

\oldpage[II]{78}

homomorphisms of quasi-coherent $\mathcal{O}_Y$-modules to the surjective
homomorphisms. Indeed, we must show that if
$\mathcal{E}\to\mathcal{E}'$ is a surjective $\mathcal{O}_Y$-homomorphism,
the diagram
\[
  \xymatrix{
    \mathbf{P}(\mathcal{E}')\times_Y\mathbf{P}(\mathcal{F})
      \ar[r]^{j\times1} \ar[d]_{\varsigma}
    & \mathbf{P}(\mathcal{E})\times_Y\mathbf{P}(\mathcal{F})
      \ar[d]^{\varsigma}
    \\\mathbf{P}(\mathcal{E}'\otimes\mathcal{F}) \ar[r]
    & \mathbf{P}(\mathcal{E}\otimes\mathcal{F})
  }
\]
is commutative, $j$ denoting the canonical closed immersion
$\mathbf{P}(\mathcal{E}')\to\mathbf{P}(\mathcal{E})$. Set
$P_1'=\mathbf{P}(\mathcal{E}')$ and otherwise keep the notation of
(\hyperref[II.4.3.1]{4.3.1}); $j\times1$ is a closed immersion
(\hyperref[I.4.3.1]{I, 4.3.1}) and we have, up to isomorphisms,
\[
  (j\times1)^*(\mathcal{O}_{P_1}(1)\otimes\mathcal{O}_{P_2}(1))
  =j^*(\mathcal{O}_{P_1}(1))\otimes\mathcal{O}_{P_2}(1)
  =\mathcal{O}_{P_1'}(1)\otimes\mathcal{O}_{P_2}(1)
\]
by virtue of (\hyperref[II.4.1.2.1]{4.1.2.1}) and of
(\hyperref[I.9.1.5]{I, 9.1.5}); our assertion therefore follows from
(\hyperref[II.4.2.8]{4.2.8}) and
(\hyperref[II.4.2.9]{4.2.9}).
\end{env}

\begin{env}[4.3.5]
\label{II.4.3.5}
With the notation of (\hyperref[II.4.3.1]{4.3.1}), let
$\psi:Y'\to Y$ be a morphism, and set
$\mathcal{E}'=\psi^*(\mathcal{E})$, $\mathcal{F}'=\psi^*(\mathcal{F})$;
then the Segre morphism
$\mathbf{P}(\mathcal{E}')\times\mathbf{P}(\mathcal{F}')
\to\mathbf{P}(\mathcal{E}'\otimes\mathcal{F}')$ is identified with
$\varsigma_{(Y')}$. Indeed, keeping the notation of
(\hyperref[II.4.3.1]{4.3.1}), let us further set
$P_1'=\mathbf{P}(\mathcal{E}')$, $P_2'=\mathbf{P}(\mathcal{F}')$; we know
(\hyperref[II.4.1.3.1]{4.1.3.1}) that $P_i'$ is identified with $(P_i)_{(Y')}$
($i=1,2$), hence the structural morphism $P_1'\times P_2'\to Y'$ is identified
with $r_{(Y')}$. On the other hand,
$\mathcal{E}'\otimes\mathcal{F}'$ is identified with
$\psi^*(\mathcal{E}\otimes\mathcal{F})$, hence
$\mathbf{P}(\mathcal{E}'\otimes\mathcal{F}')$ is identified with
$(\mathbf{P}(\mathcal{E}\otimes\mathcal{F}))_{(Y')}$
(\hyperref[II.4.1.3.1]{4.1.3.1}). Finally,
$\mathcal{O}_{P_1'}(1)\otimes_{Y'}\mathcal{O}_{P_2'}(1)=\mathcal{L}'$
is identified with $\mathcal{L}\otimes_Y\mathcal{O}_{Y'}$, by virtue of
(\hyperref[II.4.1.3.2]{4.1.3.2}) and of
(\hyperref[I.9.1.11]{I, 9.1.11}). The canonical homomorphism
$(r_{(Y')})^*(\mathcal{E}'\otimes\mathcal{F}')\to\mathcal{L}'$ is then
identified with $s_{(Y')}$, and our assertion follows from
(\hyperref[II.4.2.10]{4.2.10}).
\end{env}

\begin{remark}[4.3.6]
\label{II.4.3.6}
The sum prescheme of $\mathbf{P}(\mathcal{E})$ and
$\mathbf{P}(\mathcal{F})$ is likewise canonically isomorphic to a
closed sub-prescheme of $\mathbf{P}(\mathcal{E}\oplus\mathcal{F})$.
Indeed, the surjective homomorphisms
$\mathcal{E}\oplus\mathcal{F}\to\mathcal{E}$,
$\mathcal{E}\oplus\mathcal{F}\to\mathcal{F}$ correspond to closed
immersions
$\mathbf{P}(\mathcal{E})\to\mathbf{P}(\mathcal{E}\oplus\mathcal{F})$,
$\mathbf{P}(\mathcal{F})\to\mathbf{P}(\mathcal{E}\oplus\mathcal{F})$;
everything comes down to seeing that the underlying spaces of the closed
sub-preschemes of $\mathbf{P}(\mathcal{E}\oplus\mathcal{F})$ thus obtained have
no point in common. The question being local on $Y$, we may adopt the notation of
(\hyperref[II.4.3.1]{4.3.1}); now $\mathbf{S}_n(E)$ and
$\mathbf{S}_n(F)$ are identified with submodules of
$\mathbf{S}_n(E\oplus F)$ whose intersection is reduced to $0$; and if
$\mathfrak{p}$ is a graded prime ideal of $\mathbf{S}(E)$ such that
$\mathfrak{p}\cap\mathbf{S}_n(E)\neq\mathbf{S}_n(E)$ for every $n\geq0$,
there corresponds to it in $\mathbf{S}(E\oplus F)$ a graded prime ideal whose
trace on $\mathbf{S}_n(E)$ is
$\mathfrak{p}\cap\mathbf{S}_n(E)$, but which contains $\mathbf{S}_+(F)$,
as we see at once; two points of
$\operatorname{Proj}(\mathbf{S}(E))$ and
$\operatorname{Proj}(\mathbf{S}(F))$ respectively therefore cannot have
the same image in $\operatorname{Proj}(\mathbf{S}(E\oplus F))$.
\end{remark}


\subsection{Immersions into projective bundles. Very ample sheaves}
\label{subsection:II.4.4}

\begin{proposition}[4.4.1]
\label{II.4.4.1}
Let $Y$ be a quasi-compact scheme, or a prescheme whose underlying space
is Noetherian, $q:X\to Y$ a morphism of finite type, and $\mathcal{L}$ an
invertible $\mathcal{O}_X$-module.
\begin{enumerate}
  \item[\rm{(i)}] Let $\mathcal{S}$ be a quasi-coherent $\mathcal{O}_Y$-algebra
  graded in positive degrees, and
  $\psi:q^*(\mathcal{S})\to\bigoplus_{n\geq0}\mathcal{L}^{\otimes n}$ a
  homomorphism of graded algebras. For $r_{\mathcal{L},\psi}$ to be
  everywhere defined and to be an immersion, it is necessary and

\oldpage[II]{79}

  sufficient that there exist an integer $n>0$ and a quasi-coherent
  sub-$\mathcal{O}_Y$-module of finite type $\mathcal{E}$ of $\mathcal{S}_n$
  such that the homomorphism
  \[
    \psi'=\psi_n\circ q^*(j):q^*(\mathcal{E})
      \longrightarrow\mathcal{L}^{\otimes n}=\mathcal{L}'
  \]
  ($j$ being the injection $\mathcal{E}\to\mathcal{S}_n$) is surjective and
  the morphism
  $r_{\mathcal{L}',\psi'}:X\to\mathbf{P}(\mathcal{E})$ is an immersion.
  \item[\rm{(ii)}] Let $\mathcal{F}$ be a quasi-coherent
  $\mathcal{O}_Y$-module and $\varphi:q^*(\mathcal{F})\to\mathcal{L}$ a
  surjective homomorphism. For the morphism $r_{\mathcal{L},\varphi}$ to be an
  immersion $X\to\mathbf{P}(\mathcal{F})$, it is necessary and sufficient that
  there exist a quasi-coherent sub-$\mathcal{O}_Y$-module of finite type
  $\mathcal{E}$ of $\mathcal{F}$ such that the homomorphism
  $\varphi'=\varphi\circ q^*(j):q^*(\mathcal{E})\to\mathcal{L}$ (where
  $j:\mathcal{E}\to\mathcal{F}$ is the canonical injection) is surjective and
  $r_{\mathcal{L},\varphi'}:X\to\mathbf{P}(\mathcal{E})$ is an
  immersion.
\end{enumerate}
\end{proposition}

\begin{proof}
\begin{enumerate}
  \item[\rm{(i)}] The fact that $r_{\mathcal{L},\psi}$ is everywhere defined
  and is an immersion is equivalent, by
  (\hyperref[II.3.8.5]{3.8.5}), to the existence of $n>0$ and of
  $\mathcal{E}$ such that, if $\mathcal{S}'$ is the subalgebra of
  $\mathcal{S}$ generated by $\mathcal{E}$, the homomorphism
  $q^*(\mathcal{E})\to\mathcal{L}^{\otimes n}$ is surjective and
  $r_{\mathcal{L}',\psi'}:X\to\operatorname{Proj}(\mathcal{S}')$ is everywhere
  defined and is an immersion. On the other hand, we have a canonical
  surjective homomorphism $\mathbf{S}(\mathcal{E})\to\mathcal{S}'$, to which
  there corresponds a closed immersion
  $\operatorname{Proj}(\mathcal{S}')\to\mathbf{P}(\mathcal{E})$
  (\hyperref[II.3.6.2]{3.6.2}); whence the conclusion.
  \item[\rm{(ii)}] Since $\mathcal{F}$ is the inductive limit of its
  quasi-coherent submodules of finite type $\mathcal{E}_\lambda$
  (\hyperref[I.9.4.9]{I, 9.4.9}), $\mathbf{S}(\mathcal{F})$ is the inductive
  limit of the $\mathbf{S}(\mathcal{E}_\lambda)$; the conclusion then follows
  from (\hyperref[II.3.8.4]{3.8.4}), observing that, in the
  proof of (\hyperref[II.3.8.4]{3.8.4}), we can take all the
  $n_i$ equal to $1$: indeed ($Y$ being assumed affine), if
  $r=r_{\mathcal{L},\varphi}$ is an immersion, $r(X)$ is a quasi-compact
  subspace of $\mathbf{P}(\mathcal{F})$ which can be covered by finitely
  many open subsets of $\mathbf{P}(\mathcal{F})$ of the form $D_+(f)$ with
  $f\in F$, such that $D_+(f)\cap r(X)$ is closed.
\end{enumerate}
\end{proof}

\begin{definition}[4.4.2]
\label{II.4.4.2}
Let $Y$ be a prescheme and $q:X\to Y$ a morphism. We say that an invertible
$\mathcal{O}_X$-module $\mathcal{L}$ is \emph{very ample for
$q$} (or \emph{very ample for $Y$}, or simply \emph{very ample} if
no confusion results) if there exist a quasi-coherent $\mathcal{O}_Y$-module
$\mathcal{E}$ and a $Y$-immersion $i$ of $X$ into
$P=\mathbf{P}(\mathcal{E})$ such that $\mathcal{L}$ is isomorphic to
$i^*(\mathcal{O}_P(1))$.
\end{definition}

It amounts to the same (\hyperref[II.4.2.3]{4.2.3}) to say that there exist a
quasi-coherent $\mathcal{O}_Y$-module $\mathcal{E}$ and a surjective
homomorphism $\varphi:q^*(\mathcal{E})\to\mathcal{L}$ such that
$r_{\mathcal{L},\varphi}:X\to\mathbf{P}(\mathcal{E})$ is an immersion.

Note that the existence of an $\mathcal{O}_X$-module very ample relative
to $Y$ implies that $q$ is separated
((\hyperref[II.3.1.3]{3.1.3}) and
(\hyperref[I.5.5.1]{I, 5.5.1, (i) and (ii)})).

\begin{corollary}[4.4.3]
\label{II.4.4.3}
Suppose that there exist a quasi-coherent graded $\mathcal{O}_Y$-algebra
$\mathcal{S}$ generated by $\mathcal{S}_1$ and a $Y$-immersion
$i:X\to P=\operatorname{Proj}(\mathcal{S})$ such that $\mathcal{L}$ is
isomorphic to $i^*(\mathcal{O}_P(1))$; then $\mathcal{L}$ is very ample
relative to $q$.
\end{corollary}

\begin{proof}
Indeed, if $\mathcal{F}=\mathcal{S}_1$, the canonical homomorphism
$\mathbf{S}(\mathcal{F})\to\mathcal{S}$ is surjective; hence, by composing
the corresponding closed immersion
$\operatorname{Proj}(\mathcal{S})\to\mathbf{P}(\mathcal{F})$
(\hyperref[II.3.6.2]{3.6.2}) and the immersion $i$, we obtain an immersion
$j:X\to\mathbf{P}(\mathcal{F})=P'$ such that $\mathcal{L}$ is isomorphic to
$j^*(\mathcal{O}_{P'}(1))$ (\hyperref[II.3.6.3]{3.6.3}).
\end{proof}

\begin{proposition}[4.4.4]
\label{II.4.4.4}
Let $q:X\to Y$ be a quasi-compact morphism and $\mathcal{L}$ an invertible
$\mathcal{O}_X$-module. The following properties are
equivalent:
\begin{enumerate}
  \item[\rm{a)}] $\mathcal{L}$ is very ample relative to $q$.
  \item[\rm{b)}] $q_*(\mathcal{L})$ is quasi-coherent, the canonical
  homomorphism $\sigma:q^*(q_*(\mathcal{L}))\to\mathcal{L}$ is surjective, and
  the morphism
  $r_{\mathcal{L},\sigma}:X\to\mathbf{P}(q_*(\mathcal{L}))$ is an immersion.
\end{enumerate}
\end{proposition}

\begin{proof}
Since $q$ is quasi-compact, we know that $q_*(\mathcal{L})$ is quasi-coherent
if $q$ is separated (\hyperref[I.9.2.2]{I, 9.2.2}).

\oldpage[II]{80}

We know (\hyperref[II.3.4.7]{3.4.7}) that the existence of a surjective
homomorphism $\varphi:q^*(\mathcal{E})\to\mathcal{L}$ ($\mathcal{E}$ being a
quasi-coherent $\mathcal{O}_Y$-module) implies that $\sigma$ is surjective;
moreover, to the factorization
$q^*(\mathcal{E})\to q^*(q_*(\mathcal{L}))\xrightarrow{\sigma}\mathcal{L}$
of $\varphi$ there corresponds canonically a factorization
\[
  q^*(\mathbf{S}(\mathcal{E}))
    \longrightarrow q^*(\mathbf{S}(q_*(\mathcal{L})))
    \longrightarrow\bigoplus_{n\geq0}\mathcal{L}^{\otimes n}
\]
hence (\hyperref[II.3.8.3]{3.8.3}) the hypothesis that
$r_{\mathcal{L},\varphi}$ is an immersion implies that the same holds for
$j=r_{\mathcal{L},\sigma}$; moreover
(\hyperref[II.4.2.4]{4.2.4}), $\mathcal{L}$ is isomorphic to
$j^*(\mathcal{O}_{P'}(1))$, setting
$P'=\mathbf{P}(q_*(\mathcal{L}))$. We therefore see that a) and b) are
equivalent.
\end{proof}

\begin{corollary}[4.4.5]
\label{II.4.4.5}
Suppose $q$ quasi-compact. For $\mathcal{L}$ to be very ample
relative to $Y$, it is necessary and sufficient that there exist an open
covering $(U_\alpha)$ of $Y$ such that
$\mathcal{L}|q^{-1}(U_\alpha)$ is very ample relative to $U_\alpha$ for
every $\alpha$.
\end{corollary}

\begin{proof}
Indeed, condition b) of (\hyperref[II.4.4.4]{4.4.4}) is local on
$Y$.
\end{proof}

\begin{proposition}[4.4.6]
\label{II.4.4.6}
Let $Y$ be a quasi-compact scheme, or a prescheme whose underlying space
is Noetherian, $q:X\to Y$ a morphism of finite type, and $\mathcal{L}$ an
invertible $\mathcal{O}_X$-module. Conditions a) and b) of
(\hyperref[II.4.4.4]{4.4.4}) are then equivalent to the following:
\begin{enumerate}
  \item[\rm{a')}] There exist a quasi-coherent $\mathcal{O}_Y$-module of finite
  type $\mathcal{E}$ and a surjective homomorphism
  $\varphi:q^*(\mathcal{E})\to\mathcal{L}$ such that
  $r_{\mathcal{L},\varphi}$ is an immersion.
  \item[\rm{b')}] There exists a quasi-coherent sub-$\mathcal{O}_Y$-module of
  finite type $\mathcal{E}$ of $q_*(\mathcal{L})$ having the properties stated
  in a').
\end{enumerate}
\end{proposition}

\begin{proof}
It is clear that a') or b') implies a); on the other hand, a) implies a') by
virtue of (\hyperref[II.4.4.1]{4.4.1}), and likewise b) implies b').
\end{proof}

\begin{corollary}[4.4.7]
\label{II.4.4.7}
Suppose that $Y$ is a quasi-compact scheme or a Noetherian prescheme. If
$\mathcal{L}$ is very ample for $q$, there exist a quasi-coherent graded
$\mathcal{O}_Y$-algebra $\mathcal{S}$ such that $\mathcal{S}_1$ is of finite
type and generates $\mathcal{S}$, and a dominant open $Y$-immersion
$i:X\to P=\operatorname{Proj}(\mathcal{S})$ such that $\mathcal{L}$ is
isomorphic to $i^*(\mathcal{O}_P(1))$.
\end{corollary}

\begin{proof}
Indeed, condition b') of (\hyperref[II.4.4.6]{4.4.6}) is satisfied;
the structural morphism $p:\mathbf{P}(\mathcal{E})=P'\to Y$ is then separated
and of finite type (\hyperref[II.3.1.3]{3.1.3}), so $P'$ is a
quasi-compact scheme (resp. a Noetherian prescheme) if $Y$ is a quasi-compact
scheme (resp. a Noetherian prescheme). Let $Z$ be the closure
(\hyperref[I.9.5.10]{I, 9.5.10}) of the sub-prescheme $X'$ of $P'$
associated to the immersion $j=r_{\mathcal{L},\varphi}$ of $X$ into $P'$; it
is clear that $j$ factors into the canonical injection $Z\to P'$ and a
dominant open immersion $i:X\to Z$. But $Z$ is identified with a prescheme
$\operatorname{Proj}(\mathcal{S})$, where $\mathcal{S}$ is a graded
$\mathcal{O}_Y$-algebra that is a quotient of $\mathbf{S}(\mathcal{E})$ by a
quasi-coherent graded sheaf of ideals
(\hyperref[II.3.6.2]{3.6.2}), and it is clear that $\mathcal{S}_1$ is of
finite type and generates $\mathcal{S}$; moreover, $\mathcal{O}_Z(1)$ is the
inverse image of $\mathcal{O}_{P'}(1)$ under the canonical injection
(\hyperref[II.3.6.3]{3.6.3}), so
$\mathcal{L}=i^*(\mathcal{O}_Z(1))$.
\end{proof}

\begin{proposition}[4.4.8]
\label{II.4.4.8}
Let $q:X\to Y$ be a morphism, $\mathcal{L}$ an $\mathcal{O}_X$-module very
ample relative to $q$, and $\mathcal{L}'$ an invertible $\mathcal{O}_X$-module
such that there exist a quasi-coherent $\mathcal{O}_Y$-module $\mathcal{E}'$
and a surjective homomorphism $q^*(\mathcal{E}')\to\mathcal{L}'$. Then
$\mathcal{L}\otimes_{\mathcal{O}_X}\mathcal{L}'$ is very ample relative
to $q$.
\end{proposition}

\begin{proof}
Indeed, the hypothesis implies the existence of a $Y$-morphism
$r':X\to P'=\mathbf{P}(\mathcal{E}')$ such that
$\mathcal{L}'=r'^*(\mathcal{O}_{P'}(1))$
(\hyperref[II.4.2.1]{4.2.1}). By hypothesis there are a quasi-coherent
$\mathcal{O}_Y$-module $\mathcal{E}$ and a

\oldpage[II]{81}

$Y$-immersion $r:X\to P=\mathbf{P}(\mathcal{E})$ such that
$\mathcal{L}=r^*(\mathcal{O}_P(1))$. Set
$Q=\mathbf{P}(\mathcal{E}\otimes\mathcal{E}')$ and consider the Segre
morphism $\varsigma:P\times_YP'\to Q$ (\hyperref[II.4.3.1]{4.3.1}). Since
$r$ is an immersion, so is
$(r,r')_Y:X\to P\times_YP'$ (\hyperref[I.5.3.14]{I, 5.3.14}); but since
$\varsigma$ is an immersion (\hyperref[II.4.3.3]{4.3.3}), so is
$r'':X\xrightarrow{(r,r')_Y}P\times_YP'\xrightarrow{\varsigma}Q$. On the other
hand (\hyperref[II.4.3.2.1]{4.3.2.1}),
$\varsigma^*(\mathcal{O}_Q(1))$ is isomorphic to
$\mathcal{O}_P(1)\otimes_Y\mathcal{O}_{P'}(1)$, so
(\hyperref[I.9.1.4]{I, 9.1.4}) $r''^*(\mathcal{O}_Q(1))$ is isomorphic to
$\mathcal{L}\otimes\mathcal{L}'$, which proves the proposition.
\end{proof}

\begin{corollary}[4.4.9]
\label{II.4.4.9}
Let $q:X\to Y$ be a morphism.
\begin{enumerate}
  \item[\rm{(i)}] Let $\mathcal{L}$ be an invertible $\mathcal{O}_X$-module
  and $\mathcal{K}$ an invertible $\mathcal{O}_Y$-module. For $\mathcal{L}$
  to be very ample relative to $q$, it is necessary and sufficient that
  $\mathcal{L}\otimes q^*(\mathcal{K})$ be so.
  \item[\rm{(ii)}] If $\mathcal{L}$ and $\mathcal{L}'$ are two
  $\mathcal{O}_X$-modules very ample relative to $q$, then so is
  $\mathcal{L}\otimes\mathcal{L}'$; in particular,
  $\mathcal{L}^{\otimes n}$ is very ample relative to $q$ for every
  $n>0$.
\end{enumerate}
\end{corollary}

\begin{proof}
(ii) is an immediate consequence of (\hyperref[II.4.4.8]{4.4.8}), as is
the necessity of condition (i); on the other hand, if
$\mathcal{L}\otimes q^*(\mathcal{K})$ is very ample, then so is
$(\mathcal{L}\otimes q^*(\mathcal{K}))\otimes q^*(\mathcal{K}^{-1})$
by what precedes, and this last $\mathcal{O}_X$-module is isomorphic to
$\mathcal{L}$ (\hyperref[0.5.4.3]{0, 5.4.3} and
\hyperref[0.5.4.5]{5.4.5}).
\end{proof}

\begin{proposition}[4.4.10]
\label{II.4.4.10}
\begin{enumerate}
  \item[\rm{(i)}] For every prescheme $Y$, every invertible
  $\mathcal{O}_Y$-module $\mathcal{L}$ is very ample relative to the identity
  morphism $1_Y$.
  \item[\rm{(i bis)}] Let $f:X\to Y$ be a morphism and $j:X'\to X$ an
  immersion. If $\mathcal{L}$ is an $\mathcal{O}_X$-module very ample
  relative to $f$, then $j^*(\mathcal{L})$ is very ample relative to
  $f\circ j$.
  \item[\rm{(ii)}] Let $Z$ be a quasi-compact prescheme, $f:X\to Y$ a
  morphism of finite type, $g:Y\to Z$ a quasi-compact morphism,
  $\mathcal{L}$ an $\mathcal{O}_X$-module very ample relative to $f$, and
  $\mathcal{K}$ an $\mathcal{O}_Y$-module very ample relative to $g$.
  Then there exists an integer $n_0>0$ such that
  $\mathcal{L}\otimes f^*(\mathcal{K}^{\otimes n})$ is very ample
  relative to $g\circ f$ for every $n\geq n_0$.
  \item[\rm{(iii)}] Let $f:X\to Y$ and $g:Y'\to Y$ be two morphisms, and set
  $X'=X_{(Y')}$. If $\mathcal{L}$ is an $\mathcal{O}_X$-module very ample
  relative to $f$, then $\mathcal{L}'=\mathcal{L}\otimes_Y\mathcal{O}_{Y'}$
  is an $\mathcal{O}_{X'}$-module very ample relative to $f_{(Y')}$.
  \item[\rm{(iv)}] Let $f_i:X_i\to Y_i$ ($i=1,2$) be two $S$-morphisms.
  If $\mathcal{L}_i$ is an $\mathcal{O}_{X_i}$-module very ample relative
  to $f_i$ ($i=1,2$), then $\mathcal{L}_1\otimes_S\mathcal{L}_2$ is very ample
  relative to $f_1\times_Sf_2$.
  \item[\rm{(v)}] Let $f:X\to Y$ and $g:Y\to Z$ be two morphisms. If an
  $\mathcal{O}_X$-module $\mathcal{L}$ is very ample relative to
  $g\circ f$, then $\mathcal{L}$ is very ample relative to $f$.
  \item[\rm{(vi)}] Let $f:X\to Y$ be a morphism and $j$ the canonical
  injection $X_{\mathrm{red}}\to X$. If $\mathcal{L}$ is an
  $\mathcal{O}_X$-module very ample relative to $f$, then $j^*(\mathcal{L})$
  is very ample relative to $f_{\mathrm{red}}$.
\end{enumerate}
\end{proposition}

\begin{proof}
Property (i bis) follows immediately from the definition
(\hyperref[II.4.4.2]{4.4.2}), and it is immediate that (vi) is deduced
formally from (i bis) and (v) by an argument modelled on that of
(\hyperref[I.5.5.12]{I, 5.5.12}), which we leave to the reader. To
prove (v), consider, as in
(\hyperref[I.5.5.12]{I, 5.5.12}), the factorization
$X\xrightarrow{\Gamma_f}X\times_ZY\xrightarrow{p_2}Y$, noting that
$p_2=(g\circ f)\times1_Y$. It follows from the hypothesis and from (i) and (iv)
that $\mathcal{L}\otimes_{\mathcal{O}_Z}\mathcal{O}_Y$ is very ample for
$p_2$; on the other hand, we have
$\mathcal{L}=\Gamma_f^*(\mathcal{L}\otimes_{\mathcal{O}_Z}\mathcal{O}_Y)$
(\hyperref[I.9.1.4]{I, 9.1.4}), and $\Gamma_f$ is an immersion
(\hyperref[I.5.3.11]{I, 5.3.11}); we can therefore apply (i bis).

\oldpage[II]{82}

To prove (i), we apply the definition
(\hyperref[II.4.4.2]{4.4.2}) with $\mathcal{E}=\mathcal{L}$, noting
that $\mathbf{P}(\mathcal{E})$ is then identified with $Y$
(\hyperref[II.4.1.4]{4.1.4}).

Let us prove (iii). There exist a quasi-coherent $\mathcal{O}_Y$-module
$\mathcal{E}$ and a $Y$-immersion $i:X\to\mathbf{P}(\mathcal{E})=P$ such
that $\mathcal{L}=i^*(\mathcal{O}_P(1))$; then, if we set
$\mathcal{E}'=g^*(\mathcal{E})$, $\mathcal{E}'$ is a quasi-coherent
$\mathcal{O}_{Y'}$-module and we have
$P'=\mathbf{P}(\mathcal{E}')=P_{(Y')}$
(\hyperref[II.4.1.3.1]{4.1.3.1}), $i_{(Y')}$ is an immersion of
$X_{(Y')}$ into $P'$ (\hyperref[I.4.3.2]{I, 4.3.2}), and $\mathcal{L}'$ is
isomorphic to $(i_{(Y')})^*(\mathcal{O}_{P'}(1))$
(\hyperref[II.4.2.10]{4.2.10}).

To prove (iv), note that by hypothesis there is a $Y_i$-immersion
$r_i:X_i\to P_i=\mathbf{P}(\mathcal{E}_i)$, where $\mathcal{E}_i$ is a
quasi-coherent $\mathcal{O}_{Y_i}$-module, and
$\mathcal{L}_i=r_i^*(\mathcal{O}_{P_i}(1))$ ($i=1,2$);
$r_1\times_Sr_2$ is an $S$-immersion of $X_1\times_SX_2$ into
$P_1\times_SP_2$ (\hyperref[I.4.3.1]{I, 4.3.1}), and the inverse image of
$\mathcal{O}_{P_1}(1)\otimes_S\mathcal{O}_{P_2}(1)$ under this immersion is
$\mathcal{L}_1\otimes_S\mathcal{L}_2$. On the other hand, set
$T=Y_1\times_SY_2$, and let $p_1,p_2$ be the projections of $T$ to $Y_1,Y_2$
respectively. If we set
$P_i'=\mathbf{P}(p_i^*(\mathcal{E}_i))$ ($i=1,2$), we have by
(\hyperref[II.4.1.3.1]{4.1.3.1}) $P_i'=P_i\times_{Y_i}T$, whence
\[
  P_1'\times_TP_2'
  =(P_1\times_{Y_1}T)\times_T(P_2\times_{Y_2}T)
  =P_1\times_{Y_1}(T\times_{Y_2}P_2)
  =P_1\times_{Y_1}(Y_1\times_SP_2)
  =P_1\times_SP_2
\]
up to canonical isomorphism. Likewise, we have
$\mathcal{O}_{P_i'}(1)=\mathcal{O}_{P_i}(1)\otimes_{Y_i}\mathcal{O}_T$
(\hyperref[II.4.1.3.2]{4.1.3.2}), and an analogous calculation (based in
particular on (\hyperref[I.9.1.9.1]{I, 9.1.9.1} and
\hyperref[I.9.1.2]{9.1.2})) shows that, in the preceding identification,
$\mathcal{O}_{P_1'}(1)\otimes_T\mathcal{O}_{P_2'}(1)$ is identified with
$\mathcal{O}_{P_1}(1)\otimes_S\mathcal{O}_{P_2}(1)$. We can therefore regard
$r_1\times_Sr_2$ as a $T$-immersion of $X_1\times_SX_2$ into
$P_1'\times_TP_2'$, the inverse image of
$\mathcal{O}_{P_1'}(1)\otimes_T\mathcal{O}_{P_2'}(1)$ under this immersion
being $\mathcal{L}_1\otimes_S\mathcal{L}_2$. We then finish the argument
as in (\hyperref[II.4.4.8]{4.4.8}), using the Segre morphism.

It remains to prove (ii). We may first restrict to the case where $Z$ is an
affine scheme, since in general there is a finite covering $(U_i)$ of $Z$ by
affine open subsets; if the proposition is proved for
$\mathcal{K}|g^{-1}(U_i)$, $\mathcal{L}|f^{-1}(g^{-1}(U_i))$ and an integer
$n_i$, it will suffice to take for $n_0$ the largest of the $n_i$ in order to
prove the proposition for $\mathcal{K}$ and $\mathcal{L}$
(\hyperref[II.4.4.5]{4.4.5}). The hypothesis implies that $f$ and $g$ are
separated morphisms, so $X$ and $Y$ are quasi-compact schemes.

There is an immersion $r:X\to P=\mathbf{P}(\mathcal{E})$, where $\mathcal{E}$ is
a quasi-coherent $\mathcal{O}_Y$-module of finite type and
$\mathcal{L}=r^*(\mathcal{O}_P(1))$, by virtue of
(\hyperref[II.4.4.6]{4.4.6}). We will see that there exists an
$\mathcal{O}_P$-module $\mathcal{M}$ very ample relative to the composite
morphism $P\xrightarrow{h}Y\xrightarrow{g}Z$ such that $\mathcal{O}_P(1)$ is
isomorphic to $\mathcal{M}\otimes_Y\mathcal{K}^{\otimes(-m)}$ for some integer
$m$. For $n\geq m+1$,
$\mathcal{O}_P(1)\otimes_Y\mathcal{K}^{\otimes n}$ will therefore be very ample
for $Z$ by virtue of the hypothesis and of (iv) applied to the morphisms
$h:P\to Y$ and $1_Y$; since $r$ is an immersion and
$\mathcal{L}\otimes f^*(\mathcal{K}^{\otimes n})
=r^*(\mathcal{O}_P(1)\otimes_Y\mathcal{K}^{\otimes n})$, the conclusion
will follow from (i bis). To establish our assertion about $\mathcal{O}_P(1)$,
we will use the following lemma:

\begin{lemma}[4.4.10.1]
\label{II.4.4.10.1}
Let $Z$ be a quasi-compact scheme or a prescheme whose underlying space
is Noetherian, $g:Y\to Z$ a quasi-compact morphism, $\mathcal{K}$ an
invertible $\mathcal{O}_Y$-module very ample for $g$, and $\mathcal{E}$ a
quasi-coherent $\mathcal{O}_Y$-module of finite type. Then there exists an
integer $m_0$ such that, for every $m\geq m_0$, $\mathcal{E}$ is isomorphic to
a quotient of an $\mathcal{O}_Y$-module of the form
$g^*(\mathcal{F})\otimes\mathcal{K}^{\otimes(-m)}$, where $\mathcal{F}$ is a
quasi-coherent $\mathcal{O}_Z$-module of finite type (depending on $m$).
\end{lemma}

This lemma will be proved in (\hyperref[II.4.5.10.1]{4.5.10.1}); the reader
can verify that (\hyperref[II.4.4.10]{4.4.10}) is not used anywhere
in n\textsuperscript{o}~4.5.

\oldpage[II]{83}

Granting this lemma, there exists a closed immersion $j_1$ of $P$ into
\[
  P_1=\mathbf{P}(g^*(\mathcal{F})\otimes\mathcal{K}^{\otimes(-m)})
\]
such that $\mathcal{O}_P(1)$ is isomorphic to
$j_1^*(\mathcal{O}_{P_1}(1))$ (\hyperref[II.4.1.2]{4.1.2}). On the other
hand, there exists an isomorphism of $P_1$ onto
$P_2=\mathbf{P}(g^*(\mathcal{F}))$ identifying
$\mathcal{O}_{P_1}(1)$ with
$\mathcal{O}_{P_2}(1)\otimes_Y\mathcal{K}^{\otimes(-m)}$
(\hyperref[II.4.1.4]{4.1.4}); we therefore have a closed immersion
$j_2:P\to P_2$ such that $\mathcal{O}_P(1)$ is isomorphic to
$j_2^*(\mathcal{O}_{P_2}(1))\otimes_Y\mathcal{K}^{\otimes(-m)}$. Finally,
$P_2$ is identified with $P_3\times_ZY$, where $P_3=\mathbf{P}(\mathcal{F})$,
and $\mathcal{O}_{P_2}(1)$ with
$\mathcal{O}_{P_3}(1)\otimes_Z\mathcal{O}_Y$
(\hyperref[II.4.1.3]{4.1.3}). By definition,
$\mathcal{O}_{P_3}(1)$ is very ample for $Z$; since the same is true of
$\mathcal{K}$, we conclude from (iv) that
$\mathcal{O}_{P_2}(1)\otimes_Y\mathcal{K}$ is very ample for $Z$; the same is
true of
$\mathcal{M}=j_2^*(\mathcal{O}_{P_2}(1)\otimes_Y\mathcal{K})$ by virtue of
(i bis), and $\mathcal{O}_P(1)$ is isomorphic to
$\mathcal{M}\otimes_Y\mathcal{K}^{\otimes(-m-1)}$, which completes the
proof.
\end{proof}

\begin{proposition}[4.4.11]
\label{II.4.4.11}
Let $f:X\to Y$ and $f':X'\to Y$ be two morphisms, $X''$ the sum prescheme
$X\sqcup X'$, and $f''$ the morphism $X''\to Y$ that coincides with $f$ (resp.
$f'$) on $X$ (resp. $X'$). Let $\mathcal{L}$ (resp. $\mathcal{L}'$) be an
invertible $\mathcal{O}_X$-module (resp. an invertible
$\mathcal{O}_{X'}$-module), and let $\mathcal{L}''$ be the invertible
$\mathcal{O}_{X''}$-module that coincides with $\mathcal{L}$ on $X$ and with
$\mathcal{L}'$ on $X'$. For $\mathcal{L}''$ to be very ample relative to
$f''$, it is necessary and sufficient that $\mathcal{L}$ be very ample
relative to $f$ and $\mathcal{L}'$ very ample relative to $f'$.
\end{proposition}

\begin{proof}
We are immediately reduced to the case where $Y$ is affine. If $\mathcal{L}''$
is very ample, then so are $\mathcal{L}$ and $\mathcal{L}'$, by virtue of
(\hyperref[II.4.4.10]{4.4.10, (i bis)}). Conversely, if $\mathcal{L}$ and
$\mathcal{L}'$ are very ample, it follows immediately from the definition
(\hyperref[II.4.4.2]{4.4.2}) and from
(\hyperref[II.4.3.6]{4.3.6}) that $\mathcal{L}''$ is very ample.
\end{proof}


\subsection{Ample sheaves}
\label{subsection:II.4.5}

\begin{env}[4.5.1]
\label{II.4.5.1}
Given a prescheme $X$ and an invertible $\mathcal{O}_X$-module
$\mathcal{L}$, we will set, for every $\mathcal{O}_X$-module $\mathcal{F}$
(when no confusion about $\mathcal{L}$ is possible),
$\mathcal{F}(n)=\mathcal{F}\otimes_{\mathcal{O}_X}\mathcal{L}^{\otimes n}$
($n\in\mathbf{Z}$); on the other hand, we will set
\[
  S=\bigoplus_{n\geq0}\Gamma(X,\mathcal{L}^{\otimes n})
\]
(a graded subring of the ring $\Gamma_\bullet(\mathcal{L})$ defined in
(\hyperref[0.5.4.6]{0, 5.4.6})). If we regard $X$ as a
$\mathbf{Z}$-prescheme and denote by $p$ the structural morphism
$X\to\operatorname{Spec}(\mathbf{Z})$, there is a bijective correspondence
between the homomorphisms
\[
  p^*(\widetilde{S})\longrightarrow
  \bigoplus_{n\geq0}\mathcal{L}^{\otimes n}
\]
of graded $\mathcal{O}_X$-algebras and the endomorphisms of the graded ring
$S$ (\hyperref[I.2.2.5]{I, 2.2.5}); the homomorphism
\[
  \varepsilon:p^*(\widetilde{S})\longrightarrow
  \bigoplus_{n\geq0}\mathcal{L}^{\otimes n}
\]
that thus corresponds to the identity automorphism of $S$ will be called
canonical. To it there corresponds (\hyperref[II.3.7.1]{3.7.1}) a morphism
$G(\varepsilon)\to\operatorname{Proj}(S)$, which will also be called canonical.
\end{env}

\begin{theorem}[4.5.2]
\label{II.4.5.2}
Let $X$ be a quasi-compact scheme or a prescheme whose underlying space
is Noetherian, $\mathcal{L}$ an invertible $\mathcal{O}_X$-module, and $S$
the graded ring
$\bigoplus_{n\geq0}\Gamma(X,\mathcal{L}^{\otimes n})$. The following
conditions are equivalent:
\begin{enumerate}
  \item[\rm{a)}] As $f$ runs through the set of homogeneous elements of
  $S_+$, the $X_f$ form a basis of the topology of $X$.
  \item[\rm{a')}] As $f$ runs through the set of homogeneous elements of
  $S_+$, those $X_f$ that are affine form a covering of $X$.
  \item[\rm{b)}] The canonical morphism
  $G(\varepsilon)\to\operatorname{Proj}(S)$
  (\hyperref[II.4.5.1]{4.5.1}) is everywhere defined and is a dominant
  open immersion.

\oldpage[II]{84}

  \item[\rm{b')}] The canonical morphism
  $G(\varepsilon)\to\operatorname{Proj}(S)$ is everywhere defined and is a
  homeomorphism of the underlying space of $X$ onto a subspace of
  $\operatorname{Proj}(S)$.
  \item[\rm{c)}] For every quasi-coherent $\mathcal{O}_X$-module
  $\mathcal{F}$, if we denote by $\mathcal{F}_n$ the
  sub-$\mathcal{O}_X$-module of $\mathcal{F}(n)$ generated by the sections of
  $\mathcal{F}(n)$ over $X$, then $\mathcal{F}$ is the sum of the
  sub-$\mathcal{O}_X$-modules $\mathcal{F}_n(-n)$ for the integers $n>0$.
  \item[\rm{c')}] Property c) holds for every quasi-coherent sheaf
  of ideals of $\mathcal{O}_X$.
\end{enumerate}
Moreover, if $(f_\alpha)$ is a family of homogeneous elements of $S_+$
such that the $X_{f_\alpha}$ are affine, then the restriction to
$\bigcup_\alpha X_{f_\alpha}$ of the canonical morphism
$X\to\operatorname{Proj}(S)$ is an isomorphism of
$\bigcup_\alpha X_{f_\alpha}$ onto
$\bigcup_\alpha(\operatorname{Proj}(S))_{f_\alpha}$.
\end{theorem}

\begin{proof}
It is clear that b) implies b'), and b') implies a) by virtue of formula
(\hyperref[II.3.7.3.1]{3.7.3.1}) (taking into account that
$\varepsilon^\flat$ is the identity). Condition a) implies a'), since every
$x\in X$ has an affine neighbourhood $U$ such that
$\mathcal{L}|U$ is isomorphic to $\mathcal{O}_X|U$; if
$f\in\Gamma(X,\mathcal{L}^{\otimes n})$ is such that $x\in X_f\subset U$,
then $X_f$ is also the set of $x'\in U$ such that $(f|U)(x')\neq0$, and it is
consequently an affine open subset (\hyperref[I.1.3.6]{I, 1.3.6}). To prove
that a') implies b), it suffices to prove the last assertion of the statement
and to verify, moreover, that if $X=\bigcup_\alpha X_{f_\alpha}$, condition
(iv) of (\hyperref[II.3.8.2]{3.8.2}) is satisfied. This last point
follows immediately from (\hyperref[I.9.3.1]{I, 9.3.1, (i)}). As for the
last assertion of (\hyperref[II.4.5.2]{4.5.2}), since
$X_{f_\alpha}$ is the inverse image of
$(\operatorname{Proj}(S))_{f_\alpha}$ under
$G(\varepsilon)\to\operatorname{Proj}(S)$, it suffices to apply
(\hyperref[I.9.3.2]{I, 9.3.2}). Hence a), a'), b), b') are equivalent.

To show that a') implies c), note that if $X_f$ is affine (with
$f\in S_k$), $\mathcal{F}|X_f$ is generated by its sections over
$X_f$ (\hyperref[I.1.3.9]{I, 1.3.9}); on the other hand
(\hyperref[I.9.3.1]{I, 9.3.1, (ii)}), such a section $s$ is of the form
$(t|X_f)\otimes(f|X_f)^{-m}$ where
$t\in\Gamma(X,\mathcal{F}(km))$; by definition, $t$ is also a section
of $\mathcal{F}_{km}$, so $s$ is indeed a section of
$\mathcal{F}_{km}(-km)$ over $X_f$, which proves c). It is clear
that c) implies c'), and it remains to show that c') implies a). Now let $U$
be an open neighbourhood of $x\in X$, and let $\mathcal{J}$ be a quasi-coherent
sheaf of ideals of $\mathcal{O}_X$ defining a closed sub-prescheme of $X$
having $X-U$ as underlying space
(\hyperref[I.5.2.1]{I, 5.2.1}). Hypothesis c') implies that there exist an
integer $n>0$ and a section $f$ of $\mathcal{J}(n)$ over $X$ such
that $f(x)\neq0$. But we evidently have $f\in S_n$, and $x\in X_f\subset U$,
which proves a).
\end{proof}

When $X$ is a prescheme whose underlying space is Noetherian, the
equivalent conditions of (\hyperref[II.4.5.2]{4.5.2}) imply that
$X$ is a scheme, since it is isomorphic to a sub-prescheme of the scheme
$P=\operatorname{Proj}(S)$ by virtue of
(\hyperref[II.4.5.2]{4.5.2, b)}).

\begin{definition}[4.5.3]
\label{II.4.5.3}
We say that an invertible $\mathcal{O}_X$-module $\mathcal{L}$ is ample if
$X$ is a quasi-compact scheme and if the equivalent conditions of
(\hyperref[II.4.5.2]{4.5.2}) are satisfied.
\end{definition}

It follows evidently from criterion (\hyperref[II.4.5.2]{4.5.2, a)}) that if
$\mathcal{L}$ is an ample $\mathcal{O}_X$-module, then, for every quasi-compact open subset
$U$ of $X$, $\mathcal{L}|U$ is an ample $(\mathcal{O}_X|U)$-module.

It follows from the proof of (\hyperref[II.4.5.2]{4.5.2}) that the
affine $X_f$ already form a basis of the topology of $X$. Moreover:

\begin{corollary}[4.5.4]
\label{II.4.5.4}
Let $\mathcal{L}$ be an ample $\mathcal{O}_X$-module. For every finite subspace
$Z$ of $X$ and every neighbourhood $U$ of $Z$, there exist an integer $n$ and an
$f\in\Gamma(X,\mathcal{L}^{\otimes n})$ such that $X_f$ is an affine
neighbourhood of $Z$ contained in $U$.
\end{corollary}

\begin{proof}
\oldpage[II]{85}

By virtue of (\hyperref[II.4.5.2]{4.5.2, b)}), we may restrict ourselves to proving
that for every finite subset $Z'$ of $\operatorname{Proj}(S)$ and every
open neighbourhood $U$ of $Z'$, there exists a homogeneous element $f\in S_+$ such
that $Z'\subset(\operatorname{Proj}(S))_f\subset U$
(\hyperref[II.2.4.1]{2.4.1}). Now, by definition, the closed subset $Y$,
the complement of $U$ in $\operatorname{Proj}(S)$, is of the form
$V_+(\mathfrak{I})$, where $\mathfrak{I}$ is a graded ideal of $S$ not
containing $S_+$ (\hyperref[II.2.3.2]{2.3.2}); on the other hand, the
points of $Z'$ are by definition graded prime ideals
$\mathfrak{p}_i$ of $S_+$ not containing $\mathfrak{I}$
(\hyperref[II.2.3.1]{2.3.1}). There therefore exists an element
$f\in\mathfrak{I}$ belonging to none of the $\mathfrak{p}_i$ (Bourbaki,
\emph{Alg. comm.}, chap.~II, \S~1, n\textsuperscript{o}~1, prop.~2), and since
the $\mathfrak{p}_i$ are graded, the argument given \emph{loc. cit.}
shows that we may even assume $f$ homogeneous; this element then answers the
question.
\end{proof}

\begin{proposition}[4.5.5]
\label{II.4.5.5}
Suppose that $X$ is a quasi-compact scheme, or a prescheme whose underlying
space is Noetherian. Conditions a) to c') of
(\hyperref[II.4.5.2]{4.5.2}) are also equivalent to the following:
\begin{enumerate}
  \item[\rm{d)}] For every quasi-coherent $\mathcal{O}_X$-module
  $\mathcal{F}$ of finite type, there exists an integer $n_0$ such that, for every
  $n\geq n_0$, $\mathcal{F}(n)$ is generated by its sections over
  $X$.
  \item[\rm{d')}] For every quasi-coherent $\mathcal{O}_X$-module
  $\mathcal{F}$ of finite type, there exist two integers $n>0$, $k>0$ such that
  $\mathcal{F}$ is isomorphic to a quotient of the $\mathcal{O}_X$-module
  $\mathcal{L}^{\otimes(-n)}\otimes\mathcal{O}_X^k$.
  \item[\rm{d'')}] Property d') holds for every quasi-coherent sheaf of
  ideals of finite type of $\mathcal{O}_X$.
\end{enumerate}
\end{proposition}

\begin{proof}
Since $X$ is quasi-compact, if a quasi-coherent $\mathcal{O}_X$-module of
finite type $\mathcal{F}$ is such that $\mathcal{F}(n)$ (which is of finite type)
is generated by its sections over $X$, then $\mathcal{F}(n)$ is
generated by finitely many of these sections
(\hyperref[0.5.2.3]{0, 5.2.3}); hence d) implies d'), and it is clear that
d') implies d''). Since every quasi-coherent $\mathcal{O}_X$-module
$\mathcal{G}$ is the inductive limit of its sub-$\mathcal{O}_X$-modules of
finite type (\hyperref[I.9.4.9]{I, 9.4.9}), to verify condition c')
of (\hyperref[II.4.5.2]{4.5.2}) it suffices to do so for a quasi-coherent
sheaf of ideals of $\mathcal{O}_X$ that is of finite type, and d'')
therefore implies c'). It remains to see that if $\mathcal{L}$ is ample, property
d) is satisfied. Consider a finite covering of $X$ by
$X_{f_i}$ ($f_i\in S_{n_i}$), which we may assume to be affine; replacing the
$f_i$ by suitable powers (which does not change the $X_{f_i}$),
we may assume all the $n_i$ equal to the same integer $m$. The sheaf
$\mathcal{F}|X_{f_i}$, being of finite type by hypothesis, is generated by
finitely many of its sections $h_{ij}$ over $X_{f_i}$
(\hyperref[I.1.3.13]{I, 1.3.13}); there therefore exists an integer $k_0$ such that
the section $h_{ij}\otimes f_i^{\otimes k_0}$ extends to a section of
$\mathcal{F}(k_0m)$ over $X$ for every pair $(i,j)$
(\hyperref[I.9.3.1]{I, 9.3.1}). \emph{A fortiori} the
$h_{ij}\otimes f_i^{\otimes k}$ extend to sections of
$\mathcal{F}(km)$ over $X$ for every $k\geq k_0$, and for these values
of $k$, $\mathcal{F}(km)$ is therefore generated by its sections over
$X$. For every $p$ such that $0<p<m$, $\mathcal{F}(p)$ is also of finite type,
so there is an integer $k_p$ such that
$\mathcal{F}(p)(km)=\mathcal{F}(p+km)$ is generated by its sections
over $X$ for $k\geq k_p$. Taking $n_0$ larger than all the
$k_pm$, we conclude that $\mathcal{F}(n)$ is generated by its sections
over $X$ for every $n\geq n_0$, since such an $n$ can be written $n=km+p$ with
$k\geq k_p$ and $0\leq p<m$.
\end{proof}

\begin{proposition}[4.5.6]
\label{II.4.5.6}
Let $X$ be a quasi-compact scheme and $\mathcal{L}$ an invertible
$\mathcal{O}_X$-module.
\begin{enumerate}
  \item[\rm{(i)}] Let $n$ be an integer $>0$. For $\mathcal{L}$ to be
  ample, it is necessary and sufficient that $\mathcal{L}^{\otimes n}$ be ample.
  \item[\rm{(ii)}] Let $\mathcal{L}'$ be an invertible $\mathcal{O}_X$-module
  such that, for every $x\in X$, there exist an integer $n>0$

\oldpage[II]{86}

  and a section $s'$ of $\mathcal{L}'^{\otimes n}$ over $X$ such
  that $s'(x)\neq0$. Then, if $\mathcal{L}$ is ample, so is
  $\mathcal{L}\otimes\mathcal{L}'$.
\end{enumerate}
\end{proposition}

\begin{proof}
Property (i) is an evident consequence of criterion a) of
(\hyperref[II.4.5.2]{4.5.2}), since $X_{f^{\otimes n}}=X_f$. On the other
hand, if $\mathcal{L}$ is ample, then for every $x\in X$ and every neighbourhood $U$ of
$x$ there exist $m>0$ and $f\in\Gamma(X,\mathcal{L}^{\otimes m})$ such that
$x\in X_f\subset U$ (\hyperref[II.4.5.2]{4.5.2, a)}); if
$f'\in\Gamma(X,\mathcal{L}'^{\otimes n})$ is such that $f'(x)\neq0$, we will have
$s(x)\neq0$ for
$s=f^{\otimes n}\otimes f'^{\otimes m}\in
\Gamma(X,(\mathcal{L}\otimes\mathcal{L}')^{\otimes mn})$, hence
$x\in X_s\subset X_f\subset U$, which proves that
$\mathcal{L}\otimes\mathcal{L}'$ is ample
(\hyperref[II.4.5.2]{4.5.2, a)}).
\end{proof}

\begin{corollary}[4.5.7]
\label{II.4.5.7}
The tensor product of two ample $\mathcal{O}_X$-modules is ample.
\end{corollary}

\begin{corollary}[4.5.8]
\label{II.4.5.8}
Let $\mathcal{L}$ be an ample $\mathcal{O}_X$-module and $\mathcal{L}'$ an
invertible $\mathcal{O}_X$-module; then there exists an integer $n_0>0$ such
that $\mathcal{L}^{\otimes n}\otimes\mathcal{L}'$ is ample and generated by
its sections over $X$ for $n\geq n_0$.
\end{corollary}

\begin{proof}
Indeed, it follows from (\hyperref[II.4.5.5]{4.5.5}) that there exists an integer
$m_0$ such that $\mathcal{L}^{\otimes m}\otimes\mathcal{L}'$ is generated by
its sections over $X$ for $m\geq m_0$; by
(\hyperref[II.4.5.6]{4.5.6}) we can therefore take $n_0=m_0+1$.
\end{proof}

\begin{remark}[4.5.9]
\label{II.4.5.9}
Let $P=\mathrm{H}^1(X,\mathcal{O}_X^*)$ be the group of classes of
invertible $\mathcal{O}_X$-modules (\hyperref[0.5.4.7]{0, 5.4.7}), and
let $P^+$ be the subset of $P$ formed by the classes of ample sheaves. Suppose
that $P^+$ is non-empty. Then it follows from
(\hyperref[II.4.5.7]{4.5.7}) and
(\hyperref[II.4.5.8]{4.5.8}) that we have
\[
  P^++P^+\subset P^+
  \qquad\text{and}\qquad
  P^+-P^+=P
\]
in other words, $P^+\cup\{0\}$ is the set of positive elements of $P$
for a preorder structure on $P$ compatible with its group
structure, and which is even Archimedean by virtue of
(\hyperref[II.4.5.8]{4.5.8}). This is why one sometimes says ``positive
sheaf'' instead of ample sheaf, and ``negative sheaf'' for the inverse
of such a sheaf (a terminology which we will not follow).
\end{remark}

\begin{proposition}[4.5.10]
\label{II.4.5.10}
Let $Y$ be an affine scheme, $q:X\to Y$ a separated quasi-compact morphism,
and $\mathcal{L}$ an invertible $\mathcal{O}_X$-module.
\begin{enumerate}
  \item[\rm{(i)}] If $\mathcal{L}$ is very ample relative to $q$, then
  $\mathcal{L}$ is ample.
  \item[\rm{(ii)}] Suppose moreover that the morphism $q$ is of finite type.
  Then, for $\mathcal{L}$ to be ample, it is necessary and sufficient that it have
  one of the following equivalent properties:
  \begin{enumerate}
    \item[\rm{e)}] There exists $n_0>0$ such that for every integer $n\geq n_0$,
    $\mathcal{L}^{\otimes n}$ is very ample relative to $q$.
    \item[\rm{e')}] There exists $n>0$ such that $\mathcal{L}^{\otimes n}$ is
    very ample relative to $q$.
  \end{enumerate}
\end{enumerate}
\end{proposition}

\begin{proof}
The first assertion follows from the definition
(\hyperref[II.4.4.2]{4.4.2}) of a very ample $\mathcal{O}_X$-module: if
$A$ is the ring of $Y$, there are an $A$-module $E$ and a surjective
homomorphism
\[
  \psi:q^*((\mathbf{S}(E))^{\sim})\longrightarrow
  \bigoplus_{n\geq0}\mathcal{L}^{\otimes n}
\]
such that $i=r_{\mathcal{L},\psi}$ is an everywhere defined immersion
$X\to P=\mathbf{P}(\widetilde{E})$ and such that we have
$\mathcal{L}=i^*(\mathcal{O}_P(1))$; since the $D_+(f)$ for $f$ homogeneous
in $(\mathbf{S}(E))_+$ form a basis of the topology of $P$, and since
$i^{-1}(D_+(f))=X_{\psi^\flat(f)}$ by
(\hyperref[II.3.7.3.1]{3.7.3.1}), we see that condition a) of
(\hyperref[II.4.5.2]{4.5.2}) is satisfied, so $\mathcal{L}$ is ample.

Let us now prove that if $q$ is of finite type and $\mathcal{L}$ is
ample, then it satisfies condition e). In the first place, it follows from criterion b)
of (\hyperref[II.4.5.2]{4.5.2}) and from
(\hyperref[II.4.4.1]{4.4.1, (i)}) that there exists

\oldpage[II]{87}

an integer $k_0>0$ such that $\mathcal{L}^{\otimes k_0}$ is very ample
relative to $q$. On the other hand, by virtue of
(\hyperref[II.4.5.5]{4.5.5}), there exists an integer $m_0$ such that, for
$m\geq m_0$, $\mathcal{L}^{\otimes m}$ is generated by its sections
over $X$. Set $n_0=k_0+m_0$; if $n\geq n_0$, we can write
$n=k_0+m$ with $m\geq m_0$, whence
$\mathcal{L}^{\otimes n}
=\mathcal{L}^{\otimes k_0}\otimes\mathcal{L}^{\otimes m}$. Since
$\mathcal{L}^{\otimes m}$ is generated by its sections over $X$, it
follows from criterion (\hyperref[II.4.4.8]{4.4.8}) and from
(\hyperref[II.3.4.7]{3.4.7}) that $\mathcal{L}^{\otimes n}$ is very ample
relative to $q$. Finally, it is clear that e) implies e'), and e') implies
that $\mathcal{L}$ is ample by virtue of (i) and of
(\hyperref[II.4.5.6]{4.5.6, (i)}).

\begin{env}[4.5.10.1]
\label{II.4.5.10.1}
\emph{Proof of Lemma (\hyperref[II.4.4.10.1]{4.4.10.1}).} Set
$\mathcal{E}(n)=\mathcal{E}\otimes\mathcal{K}^{\otimes n}$; since $g$ is
separated (\hyperref[II.4.4.2]{4.4.2}), the argument of
(\hyperref[II.3.4.8]{3.4.8}) applies and reduces us to seeing that
the canonical homomorphism
$g^*(g_*(\mathcal{E}(n)))\to\mathcal{E}(n)$ is surjective for $n$ sufficiently
large. Moreover, since $Z$ is quasi-compact, the argument of
(\hyperref[II.3.4.6]{3.4.6}) reduces us to proving the proposition in the case
where $Z$ is affine. Now $\mathcal{K}$ is then ample by virtue of
(\hyperref[II.4.5.10]{4.5.10, (i)}), and the conclusion follows from
(\hyperref[II.4.5.5]{4.5.5, d')}).
\end{env}
\end{proof}

\begin{corollary}[4.5.11]
\label{II.4.5.11}
Let $Y$ be an affine scheme, $q:X\to Y$ a separated morphism of finite type,
$\mathcal{L}$ an ample $\mathcal{O}_X$-module, and $\mathcal{L}'$ an
invertible $\mathcal{O}_X$-module. There exists an integer $n_0$ such that, for
$n\geq n_0$, $\mathcal{L}^{\otimes n}\otimes\mathcal{L}'$ is very ample
relative to $q$.
\end{corollary}

\begin{proof}
Indeed, there exists $m_0$ such that for $m\geq m_0$,
$\mathcal{L}^{\otimes m}\otimes\mathcal{L}'$ is generated by its sections
over $X$ (\hyperref[II.4.5.8]{4.5.8}); on the other hand, there exists
$k_0$ such that $\mathcal{L}^{\otimes k}$ is very ample relative to $q$
for $k\geq k_0$. Consequently
$\mathcal{L}^{\otimes(k+m_0)}\otimes\mathcal{L}'$ is very ample for
$k\geq k_0$ ((\hyperref[II.4.4.8]{4.4.8}) and
(\hyperref[II.3.4.7]{3.4.7})).
\end{proof}

\begin{remark}[4.5.12]
\label{II.4.5.12}
It is not known whether the hypothesis that an $\mathcal{O}_X$-module $\mathcal{L}$ is such
that $\mathcal{L}^{\otimes n}$ is very ample (relative to $q$) implies
the same conclusion for $\mathcal{L}^{\otimes(n+1)}$.
\end{remark}

\begin{proposition}[4.5.13]
\label{II.4.5.13}
Let $X$ be a quasi-compact prescheme, $Z$ a closed sub-prescheme of $X$
defined by a nilpotent quasi-coherent sheaf $\mathcal{J}$ of ideals of
$\mathcal{O}_X$, and $j$ the canonical injection $Z\to X$. For an
invertible $\mathcal{O}_X$-module $\mathcal{L}$ to be ample, it is necessary and
sufficient that $\mathcal{L}'=j^*(\mathcal{L})$ be an ample
$\mathcal{O}_Z$-module.
\end{proposition}

\begin{proof}
The condition is necessary. Indeed, for every section $f$ of
$\mathcal{L}^{\otimes n}$ over $X$, let $f'$ be its canonical image
$f\otimes1$, a section of
$\mathcal{L}'^{\otimes n}=\mathcal{L}^{\otimes n}
\otimes_{\mathcal{O}_X}(\mathcal{O}_X/\mathcal{J})$ over the space
$Z$ (which is identical to $X$); it is clear that $X_f=Z_{f'}$, so
criterion a) of (\hyperref[II.4.5.2]{4.5.2}) shows that $\mathcal{L}'$ is
ample.

To see that the condition is sufficient, note first that we can
reduce to the case where $\mathcal{J}^2=0$ by considering the (finite) sequence of
preschemes
$X_k=(X,\mathcal{O}_X/\mathcal{J}^{k+1})$, each of which is a closed
sub-prescheme of the next, defined by a sheaf of ideals of square zero. Moreover,
$X$ is a scheme, since $X_{\mathrm{red}}$ is one by hypothesis
((\hyperref[II.4.5.3]{4.5.3}) and
(\hyperref[I.5.5.1]{I, 5.5.1})). Criterion a) of
(\hyperref[II.4.5.2]{4.5.2}) shows that it will suffice to prove

\begin{lemma}[4.5.13.1]
\label{II.4.5.13.1}
Under the hypotheses of (\hyperref[II.4.5.13]{4.5.13}), suppose moreover
that $\mathcal{J}$ is of square zero; $\mathcal{L}$ being an
invertible $\mathcal{O}_X$-module, let $g$ be a section of
$\mathcal{L}'^{\otimes n}$ over $Z$ such that $Z_g$ is affine.
Then there exists an integer $m>0$ such that $g^{\otimes m}$ is the
canonical image of a section $f$ of $\mathcal{L}^{\otimes nm}$ over $X$.
\end{lemma}

We have the exact sequence of $\mathcal{O}_X$-modules
\[
  0\longrightarrow\mathcal{J}(n)\longrightarrow
  \mathcal{O}_X(n)=\mathcal{L}^{\otimes n}\longrightarrow
  \mathcal{O}_Z(n)=\mathcal{L}'^{\otimes n}\longrightarrow0
\]

\oldpage[II]{88}

since $\mathcal{F}(n)$ is an exact functor in $\mathcal{F}$; whence the
exact cohomology sequence
\[
  0\longrightarrow\Gamma(X,\mathcal{J}(n))\longrightarrow
  \Gamma(X,\mathcal{L}^{\otimes n})\longrightarrow
  \Gamma(X,\mathcal{L}'^{\otimes n})\xrightarrow{\partial}
  \mathrm{H}^1(X,\mathcal{J}(n))
\]
which in particular associates to $g$ an element
$\partial g\in\mathrm{H}^1(X,\mathcal{J}(n))$.

Note that, since $\mathcal{J}^2=0$, $\mathcal{J}$ can be regarded
as a quasi-coherent $\mathcal{O}_Z$-module, and we have, for every $k$,
$\mathcal{L}'^{\otimes k}\otimes_{\mathcal{O}_Z}\mathcal{J}(n)
=\mathcal{J}(n+k)$; for every section
$s\in\Gamma(X,\mathcal{L}'^{\otimes k})$, tensor multiplication by
$s$ is therefore a homomorphism $\mathcal{J}(n)\to\mathcal{J}(n+k)$ of
$\mathcal{O}_Z$-modules, which consequently gives a homomorphism
$\mathrm{H}^i(X,\mathcal{J}(n))\to
\mathrm{H}^i(X,\mathcal{J}(n+k))$ of cohomology groups.

This being so, we will see that we have
\begin{equation}
\label{II.4.5.13.2}
  g^{\otimes m}\otimes\partial g=0
\tag{4.5.13.2}
\end{equation}
for some sufficiently large $m>0$. Indeed, $Z_g$ is an affine open subset of $Z$, and
we therefore have $\mathrm{H}^1(Z_g,\mathcal{J}(n))=0$ when $\mathcal{J}(n)$ is
regarded as an $\mathcal{O}_Z$-module
(\hyperref[I.5.1.9.2]{I, 5.1.9.2}). In particular, if we set
$g'=g|Z_g$ and consider its image under the map
$\partial:\Gamma(Z_g,\mathcal{L}'^{\otimes n})\to
\mathrm{H}^1(Z_g,\mathcal{J}(n))$, we have $\partial g'=0$. To make
this relation explicit, observe that in dimension $1$ the cohomology of a sheaf
of abelian groups is the same as its Čech cohomology (G, II, 5.9);
to form $\partial g$, we must therefore consider a sufficiently fine open covering
$(U_\alpha)$ of $X$, which we may assume finite and formed
of affine open subsets, take for each $\alpha$ a section
$g_\alpha\in\Gamma(U_\alpha,\mathcal{L}^{\otimes n})$ whose canonical
image in $\Gamma(U_\alpha,\mathcal{L}'^{\otimes n})$ is
$g|U_\alpha$, and consider the class of the cocycle
$(g_{\alpha|\beta}-g_{\beta|\alpha})$, $g_{\alpha|\beta}$ being the
restriction of $g_\alpha$ to $U_\alpha\cap U_\beta$ (this cocycle having
values in $\mathcal{J}(n)$). We may moreover assume that $\partial g'$
is computed in the same way by means of the covering formed by the
$U_\alpha\cap Z_g$ and the restrictions $g_\alpha|(U_\alpha\cap Z_g)$
(replacing $(U_\alpha)$ by a finer covering if necessary); the
relation $\partial g'=0$ then means that there exists for each $\alpha$
a section $h_\alpha\in\Gamma(U_\alpha\cap Z_g,\mathcal{J}(n))$ such that
\[
  (g_{\alpha|\beta}-g_{\beta|\alpha})
  |(U_\alpha\cap U_\beta\cap Z_g)
  =h_{\alpha|\beta}-h_{\beta|\alpha}
\]
where we denote by $h_{\alpha|\beta}$ the restriction of $h_\alpha$ to
$U_\alpha\cap U_\beta\cap Z_g$ (G, II, 5.11). This being so, there is an integer
$m>0$ such that $g^{\otimes m}\otimes h_\alpha$ is the restriction to
$U_\alpha\cap Z_g$ of a section
$t_\alpha\in\Gamma(U_\alpha,\mathcal{J}(n+nm))$ for every $\alpha$
(\hyperref[I.9.3.1]{I, 9.3.1}); we therefore have
\[
  g^{\otimes m}\otimes(g_{\alpha|\beta}-g_{\beta|\alpha})
  =t_{\alpha|\beta}-t_{\beta|\alpha}
\]
for every pair of indices, which proves
(\hyperref[II.4.5.13.2]{4.5.13.2}).

Note, on the other hand, that if
$s\in\Gamma(X,\mathcal{O}_Z(p))$ and
$t\in\Gamma(X,\mathcal{O}_Z(q))$, we have, in the group
$\mathrm{H}^1(X,\mathcal{J}(p+q))$,
\begin{equation}
\label{II.4.5.13.3}
  \partial(s\otimes t)=(\partial s)\otimes t+s\otimes(\partial t).
\tag{4.5.13.3}
\end{equation}

Indeed, to compute the two sides, we can again consider an open
covering $(U_\alpha)$ of $X$ and, for each $\alpha$, a section
$s_\alpha\in\Gamma(U_\alpha,\mathcal{O}_X(p))$ (resp.
$t_\alpha\in\Gamma(U_\alpha,\mathcal{O}_X(q))$) whose canonical image in
$\Gamma(U_\alpha,\mathcal{O}_Z(p))$ (resp.
$\Gamma(U_\alpha,\mathcal{O}_Z(q))$ is $s|U_\alpha$ (resp.
$t|U_\alpha$); relation
(\hyperref[II.4.5.13.3]{4.5.13.3}) then follows from the relations
\[
  (s_{\alpha|\beta}\otimes t_{\alpha|\beta})
  -(s_{\beta|\alpha}\otimes t_{\beta|\alpha})
  =(s_{\alpha|\beta}-s_{\beta|\alpha})\otimes t_{\alpha|\beta}
  +s_{\beta|\alpha}\otimes(t_{\alpha|\beta}-t_{\beta|\alpha})
\]
with the same notation as above. By induction on $k$, we therefore have
\begin{equation}
\label{II.4.5.13.4}
  \partial(g^{\otimes k})=(kg^{\otimes(k-1)})\otimes(\partial g)
\tag{4.5.13.4}
\end{equation}

\oldpage[II]{89}

and we conclude from (\hyperref[II.4.5.13.2]{4.5.13.2}) and
(\hyperref[II.4.5.13.4]{4.5.13.4}) that
$\partial(g^{\otimes(m+1)})=0$; $g^{\otimes(m+1)}$ is therefore the
canonical image of a section $f$ of $\mathcal{L}^{\otimes n(m+1)}$ over
$X$, which completes the proof of (\hyperref[II.4.5.13]{4.5.13}).
\end{proof}

\begin{corollary}[4.5.14]
\label{II.4.5.14}
Let $X$ be a Noetherian scheme and $j$ the canonical injection
$X_{\mathrm{red}}\to X$. For an invertible $\mathcal{O}_X$-module
$\mathcal{L}$ to be ample, it is necessary and sufficient that $j^*(\mathcal{L})$ be
an ample $\mathcal{O}_{X_{\mathrm{red}}}$-module.
\end{corollary}

This follows from (\hyperref[I.6.1.6]{I, 6.1.6}).


\subsection{Relatively ample sheaves}
\label{II.4.6}

\begin{definition}[4.6.1]
\label{II.4.6.1}
Let $f:X\to Y$ be a quasi-compact morphism and $\mathcal{L}$ an invertible
$\mathcal{O}_X$-module. We say that $\mathcal{L}$ is ample
relative to $f$, or relative to $Y$, or $f$-ample, or $Y$-ample (or
even ample if no confusion with the notion defined in
(\hyperref[II.4.5.3]{4.5.3}) results) if there exists a covering $(U_\alpha)$
of $Y$ formed of affine open subsets such that, if we set
$X_\alpha=f^{-1}(U_\alpha)$, $\mathcal{L}|X_\alpha$ is an ample
$\mathcal{O}_{X_\alpha}$-module for every $\alpha$.
\end{definition}

The existence of an $f$-ample $\mathcal{O}_X$-module implies that $f$ is
necessarily separated ((\hyperref[II.4.5.3]{4.5.3}) and
(\hyperref[I.5.5.5]{I, 5.5.5})).

\begin{proposition}[4.6.2]
\label{II.4.6.2}
Let $f:X\to Y$ be a quasi-compact morphism and $\mathcal{L}$ an invertible
$\mathcal{O}_X$-module. If $\mathcal{L}$ is very ample
relative to $f$, it is ample relative to $f$.
\end{proposition}

This follows from the local character (on $Y$) of the notion of very
ample sheaf (\hyperref[II.4.4.5]{4.4.5}), from the definition
(\hyperref[II.4.6.1]{4.6.1}) and from criterion
(\hyperref[II.4.5.10]{4.5.10, (i)}).

\begin{proposition}[4.6.3]
\label{II.4.6.3}
Let $f:X\to Y$ be a quasi-compact morphism, $\mathcal{L}$ an invertible
$\mathcal{O}_X$-module, and let $\mathcal{S}$ be the graded
$\mathcal{O}_Y$-algebra
$\bigoplus_{n\geq0}f_*(\mathcal{L}^{\otimes n})$. The following conditions
are equivalent:
\begin{enumerate}
  \item[\rm{a)}] $\mathcal{L}$ is $f$-ample.
  \item[\rm{b)}] $\mathcal{S}$ is quasi-coherent and the canonical
  homomorphism $\sigma:f^*(\mathcal{S})\to
  \bigoplus_{n\geq0}\mathcal{L}^{\otimes n}$
  (\hyperref[0.4.4.3]{0, 4.4.3}) is such that the $Y$-morphism
  $r_{\mathcal{L},\sigma}:G(\sigma)\to\operatorname{Proj}(\mathcal{S})=P$
  is everywhere defined and is a dominant open immersion.
  \item[\rm{b')}] The morphism $f$ is separated, and the $Y$-morphism
  $r_{\mathcal{L},\sigma}$ is everywhere defined and is a homeomorphism of
  the underlying space of $X$ onto a subspace of
  $\operatorname{Proj}(\mathcal{S})$.
\end{enumerate}
Moreover, when this is so, for every $n\in\mathbf{Z}$,
the canonical homomorphism
\begin{equation}
\label{II.4.6.3.1}
  r_{\mathcal{L},\sigma}^*(\mathcal{O}_P(n))
  \longrightarrow\mathcal{L}^{\otimes n}
\tag{4.6.3.1}
\end{equation}
defined in (\hyperref[II.3.7.9.1]{3.7.9.1}) is an isomorphism.

Finally, for every quasi-coherent $\mathcal{O}_X$-module $\mathcal{F}$, if
we set
$\mathcal{M}=\bigoplus_{n\geq0}f_*(\mathcal{F}\otimes
\mathcal{L}^{\otimes n})$, the canonical homomorphism
\begin{equation}
\label{II.4.6.3.2}
  r_{\mathcal{L},\sigma}^*(\widetilde{\mathcal{M}})
  \longrightarrow\mathcal{F}
\tag{4.6.3.2}
\end{equation}
defined in (\hyperref[II.3.7.9.2]{3.7.9.2}) is an isomorphism.
\end{proposition}

We have remarked that a) implies that $f$ is separated, so $\mathcal{S}$ is
quasi-coherent (\hyperref[I.9.2.2]{I, 9.2.2, a)}). Since the fact that
$r_{\mathcal{L},\sigma}$ is an everywhere defined immersion is of a
local character on $Y$, to prove that a) implies b) we may assume
$Y$ affine and $\mathcal{L}$ ample;

\oldpage[II]{90}

the assertion then follows from (\hyperref[II.4.5.2]{4.5.2, b)}). It is
clear that b) implies b'); finally, to prove that b') implies a), it
suffices to consider a covering of $Y$ by affine open subsets
$U_\alpha$ and to apply criterion
(\hyperref[II.4.5.2]{4.5.2, b')}) to each sheaf
$\mathcal{L}|f^{-1}(U_\alpha)$.

For the last two assertions, we use the fact that
$\sigma^\flat$ is here the identity, and the explicit description of the homomorphisms
(\hyperref[II.3.7.9.1]{3.7.9.1}) and
(\hyperref[II.3.7.9.2]{3.7.9.2}); it follows immediately that
(\hyperref[II.4.6.3.1]{4.6.3.1}) is an isomorphism. As for
(\hyperref[II.4.6.3.2]{4.6.3.2}), we can reduce to the case where $Y$ is
affine, hence $\mathcal{L}$ ample; it is clear that the homomorphism
(\hyperref[II.4.6.3.2]{4.6.3.2}) is injective, and criterion
(\hyperref[II.4.5.2]{4.5.2, c)}) shows that it is surjective, whence the
conclusion.

\begin{corollary}[4.6.4]
\label{II.4.6.4}
Let $(U_\alpha)$ be an open covering of $Y$. For $\mathcal{L}$
to be ample relative to $Y$, it is necessary and sufficient that
$\mathcal{L}|f^{-1}(U_\alpha)$ be ample relative to $U_\alpha$ for
every $\alpha$.
\end{corollary}

Indeed, condition b) is local on $Y$.

\begin{corollary}[4.6.5]
\label{II.4.6.5}
Let $\mathcal{K}$ be an invertible $\mathcal{O}_Y$-module. For
$\mathcal{L}$ to be $Y$-ample, it is necessary and sufficient that
$\mathcal{L}\otimes f^*(\mathcal{K})$ be so.
\end{corollary}

This is an evident consequence of (\hyperref[II.4.6.4]{4.6.4}),
taking the $U_\alpha$ such that $\mathcal{K}|U_\alpha$ is isomorphic to
$\mathcal{O}_Y|U_\alpha$ for every $\alpha$.

\begin{corollary}[4.6.6]
\label{II.4.6.6}
Suppose $Y$ affine; for $\mathcal{L}$ to be $Y$-ample, it is necessary and
sufficient that $\mathcal{L}$ be ample.
\end{corollary}

This is an immediate consequence of the definition
(\hyperref[II.4.6.1]{4.6.1}) and of criteria
(\hyperref[II.4.6.3]{4.6.3, b)}) and
(\hyperref[II.4.5.2]{4.5.2, b)}), since here
$\operatorname{Proj}(\mathcal{S})
=\operatorname{Proj}(\Gamma(Y,\mathcal{S}))$ by definition.

\begin{corollary}[4.6.7]
\label{II.4.6.7}
Let $f:X\to Y$ be a quasi-compact morphism. Suppose that there exist a
quasi-coherent $\mathcal{O}_Y$-module $\mathcal{E}$ and a $Y$-morphism
$g:X\to P=\mathbf{P}(\mathcal{E})$ that is a homeomorphism of the underlying
space of $X$ onto a subspace of $P$; then
$\mathcal{L}=g^*(\mathcal{O}_P(1))$ is $Y$-ample.
\end{corollary}

Indeed, we may assume $Y$ affine; the corollary then follows from
criterion (\hyperref[II.4.5.2]{4.5.2, a)}), from formula
(\hyperref[II.3.7.3.1]{3.7.3.1}) and from
(\hyperref[II.4.2.3]{4.2.3}).

\begin{proposition}[4.6.8]
\label{II.4.6.8}
Let $X$ be a quasi-compact scheme or a prescheme whose underlying
space is Noetherian, and $f:X\to Y$ a separated quasi-compact morphism.
For an invertible $\mathcal{O}_X$-module $\mathcal{L}$ to be $f$-ample,
it is necessary and sufficient that one of the following equivalent conditions be
satisfied:
\begin{enumerate}
  \item[\rm{c)}] For every quasi-coherent $\mathcal{O}_X$-module
  $\mathcal{F}$ of finite type, there exists an integer $n_0>0$ such that, for
  every $n\geq n_0$, the canonical homomorphism
  $\sigma:f^*(f_*(\mathcal{F}\otimes\mathcal{L}^{\otimes n}))\to
  \mathcal{F}\otimes\mathcal{L}^{\otimes n}$ is surjective.
  \item[\rm{c')}] For every quasi-coherent sheaf of ideals of finite type
  $\mathcal{J}$ of $\mathcal{O}_X$, there exists an integer $n>0$ such that
  the canonical homomorphism
  $\sigma:f^*(f_*(\mathcal{J}\otimes\mathcal{L}^{\otimes n}))\to
  \mathcal{J}\otimes\mathcal{L}^{\otimes n}$ is surjective.
\end{enumerate}
\end{proposition}

Since $X$ is quasi-compact, so is $f(X)$, and there therefore exists
a finite covering $(U_i)$ of $f(X)$ by affine open subsets of $Y$. To
prove condition c) when $\mathcal{L}$ is $f$-ample, we can
replace $Y$ by the $U_i$ and $X$ by the $f^{-1}(U_i)$, since if we obtain
for each $i$ an integer $n_i$ such that c) holds (for $U_i$,
$f^{-1}(U_i)$ and $\mathcal{L}|f^{-1}(U_i)$) as soon as $n\geq n_i$, it
will suffice to take for $n_0$ the largest of the $n_i$ in order to obtain c) for
$Y$, $X$ and $\mathcal{L}$. Now when $Y$ is affine, condition c)
follows from (\hyperref[II.4.5.5]{4.5.5, d)}), taking into account
(\hyperref[II.4.6.6]{4.6.6}). It is trivial that c) implies c'). Finally,
to prove that c') implies that $\mathcal{L}$ is $f$-ample, we may
again restrict to the case where $Y$ is affine: indeed, every quasi-coherent
sheaf of ideals of finite type $\mathcal{J}_i$ of
$\mathcal{O}_X|f^{-1}(U_i)$ is the restriction of a quasi-coherent sheaf
of ideals of finite type of $\mathcal{O}_X$
(\hyperref[I.9.4.7]{I, 9.4.7}), and hypothesis c') implies that

\oldpage[II]{91}

$\mathcal{J}_i\otimes(\mathcal{L}^{\otimes n}|f^{-1}(U_i))$ is generated
by its sections (taking into account (\hyperref[I.9.2.2]{I, 9.2.2}) and
(\hyperref[II.3.4.7]{3.4.7})); it therefore suffices to apply criterion
(\hyperref[II.4.5.5]{4.5.5, d'')}).

\begin{proposition}[4.6.9]
\label{II.4.6.9}
Let $f:X\to Y$ be a quasi-compact morphism and $\mathcal{L}$ an invertible
$\mathcal{O}_X$-module.
\begin{enumerate}
  \item[\rm{(i)}] Let $n$ be an integer $>0$. For $\mathcal{L}$ to be
  $f$-ample, it is necessary and sufficient that $\mathcal{L}^{\otimes n}$ be
  $f$-ample.
  \item[\rm{(ii)}] Let $\mathcal{L}'$ be an invertible $\mathcal{O}_X$-module,
  and suppose that there exists an integer $n>0$ such that
  the canonical homomorphism
  $\sigma:f^*(f_*(\mathcal{L}'^{\otimes n}))\to
  \mathcal{L}'^{\otimes n}$ is surjective. Then, if $\mathcal{L}$ is
  $f$-ample, so is $\mathcal{L}\otimes\mathcal{L}'$.
\end{enumerate}
\end{proposition}

Indeed, we are immediately reduced to the case where $Y$ is affine, and the proposition
is then an immediate consequence of (\hyperref[II.4.5.6]{4.5.6}).

\begin{corollary}[4.6.10]
\label{II.4.6.10}
The tensor product of two $f$-ample $\mathcal{O}_X$-modules is
$f$-ample.
\end{corollary}

\begin{proposition}[4.6.11]
\label{II.4.6.11}
Let $Y$ be a quasi-compact prescheme, $f:X\to Y$ a morphism of finite
type, and $\mathcal{L}$ an invertible $\mathcal{O}_X$-module. For
$\mathcal{L}$ to be $f$-ample, it is necessary and sufficient that it have one of the
following equivalent properties:
\begin{enumerate}
  \item[\rm{d)}] There exists $n_0>0$ such that, for every integer $n\geq n_0$,
  $\mathcal{L}^{\otimes n}$ is very ample relative to $f$.
  \item[\rm{d')}] There exists $n>0$ such that $\mathcal{L}^{\otimes n}$ is
  very ample relative to $f$.
\end{enumerate}
\end{proposition}

If $\mathcal{L}$ is ample relative to $f$, there is a finite covering
$(U_i)$ of $Y$ by affine open subsets such that the
$\mathcal{L}|f^{-1}(U_i)$ are ample. We conclude from this
(\hyperref[II.4.5.10]{4.5.10}) that there exists an integer $n_0$ such that
$\mathcal{L}^{\otimes n}|f^{-1}(U_i)$ is very ample relative to
$f^{-1}(U_i)\to U_i$ for every $n\geq n_0$ and every $i$, so
$\mathcal{L}^{\otimes n}$ is very ample relative to $f$
(\hyperref[II.4.4.5]{4.4.5}). Conversely, d') already implies that
$\mathcal{L}^{\otimes n}$ is $f$-ample
(\hyperref[II.4.6.2]{4.6.2}), so the same is true of $\mathcal{L}$
(\hyperref[II.4.6.9]{4.6.9, (i)}).

\begin{corollary}[4.6.12]
\label{II.4.6.12}
Let $Y$ be a quasi-compact prescheme, $f:X\to Y$ a morphism of finite
type, and $\mathcal{L}$, $\mathcal{L}'$ two invertible
$\mathcal{O}_X$-modules. If $\mathcal{L}$ is $f$-ample, there exists $n_0$ such that
$\mathcal{L}^{\otimes n}\otimes\mathcal{L}'$ is very ample relative
to $f$ for every $n\geq n_0$.
\end{corollary}

We argue as in (\hyperref[II.4.6.11]{4.6.11}), using a finite
affine open covering of $Y$ and
(\hyperref[II.4.5.11]{4.5.11}).

\begin{proposition}[4.6.13]
\label{II.4.6.13}
\begin{enumerate}
  \item[\rm{(i)}] For every prescheme $Y$, every invertible
  $\mathcal{O}_Y$-module $\mathcal{L}$ is ample relative
  to the identity morphism $1_Y$.
  \item[\rm{(i bis)}] Let $f:X\to Y$ be a quasi-compact morphism and
  $j:X'\to X$ a quasi-compact morphism that is a homeomorphism of
  the underlying space of $X'$ onto a subspace of $X$. If $\mathcal{L}$
  is an $\mathcal{O}_X$-module ample relative to $f$, then
  $j^*(\mathcal{L})$ is ample relative to $f\circ j$.
  \item[\rm{(ii)}] Let $Z$ be a quasi-compact prescheme, $f:X\to Y$ and
  $g:Y\to Z$ two quasi-compact morphisms, $\mathcal{L}$ an
  $\mathcal{O}_X$-module ample relative to $f$, and $\mathcal{K}$ an
  $\mathcal{O}_Y$-module ample relative to $g$. Then there exists an
  integer $n_0>0$ such that
  $\mathcal{L}\otimes f^*(\mathcal{K}^{\otimes n})$ is ample relative
  to $g\circ f$ for every $n\geq n_0$.
  \item[\rm{(iii)}] Let $f:X\to Y$ be a quasi-compact morphism and
  $g:Y'\to Y$ a morphism, and set $X'=X_{(Y')}$. If $\mathcal{L}$ is
  an $\mathcal{O}_X$-module ample relative to $f$, then
  $\mathcal{L}'=\mathcal{L}\otimes_Y\mathcal{O}_{Y'}$ is an
  $\mathcal{O}_{X'}$-module ample relative to $f_{(Y')}$.
  \item[\rm{(iv)}] Let $f_i:X_i\to Y_i$ ($i=1,2$) be two quasi-compact
  $S$-morphisms. If $\mathcal{L}_i$ is an
  $\mathcal{O}_{X_i}$-module ample relative to $f_i$ ($i=1,2$), then
  $\mathcal{L}_1\otimes_S\mathcal{L}_2$ is ample relative to
  $f_1\times_S f_2$.
  \item[\rm{(v)}] Let $f:X\to Y$ and $g:Y\to Z$ be two morphisms such that
  $g\circ f$ is quasi-compact. If an $\mathcal{O}_X$-module
  $\mathcal{L}$ is ample relative to $g\circ f$, and if $g$ is separated
  or the underlying space of $X$ is locally Noetherian, then $\mathcal{L}$
  is ample relative to $f$.

  \oldpage[II]{92}

  \item[\rm{(vi)}] Let $f:X\to Y$ be a quasi-compact morphism and $j$
  the canonical injection $X_{\mathrm{red}}\to X$. If $\mathcal{L}$ is an
  $\mathcal{O}_X$-module ample relative to $f$, then
  $j^*(\mathcal{L})$ is ample relative to $f_{\mathrm{red}}$.
\end{enumerate}
\end{proposition}

Note first that (v) and (vi) are derived from (i), (i bis) and (iv) by the
same argument as in (\hyperref[II.4.4.10]{4.4.10}), using
(\hyperref[II.4.6.4]{4.6.4}) instead of
(\hyperref[II.4.4.5]{4.4.5}); we leave the details of the argument to the
reader. Assertion (i) is a trivial consequence of
(\hyperref[II.4.4.10]{4.4.10, (i)}) and of
(\hyperref[II.4.6.2]{4.6.2}). To prove (i bis), (iii) and (iv), we
will use the following lemma:

\begin{lemma}[4.6.13.1]
\label{II.4.6.13.1}
\begin{enumerate}
  \item[\rm{(i)}] Let $u:Z\to S$ be a morphism, $\mathcal{L}$ an invertible
  $\mathcal{O}_S$-module, $s$ a section of $\mathcal{L}$
  over $S$, and $s'$ the section of $u^*(\mathcal{L})=\mathcal{L}'$
  over $Z$ that corresponds canonically to it. Then we have
  $Z_{s'}=u^{-1}(S_s)$.
  \item[\rm{(ii)}] Let $Z$, $Z'$ be two $S$-preschemes, $p$, $p'$ the
  projections of $T=Z\times_S Z'$, $\mathcal{L}$ (resp.
  $\mathcal{L}'$) an invertible $\mathcal{O}_Z$-module (resp. an invertible
  $\mathcal{O}_{Z'}$-module), $t$ (resp. $t'$) a section of
  $\mathcal{L}$ (resp. $\mathcal{L}'$) over $Z$ (resp. $Z'$), and $s$
  (resp. $s'$) the section of $p^*(\mathcal{L})$ (resp.
  $p'^*(\mathcal{L}')$) over $Z\times_S Z'$ that corresponds to it.
  Then we have $T_{s\otimes s'}=Z_t\times_S Z'_{t'}$.
\end{enumerate}
\end{lemma}

It follows from the definitions that we can reduce to the case where all the
preschemes considered are affine. Moreover, in (i), we may assume
$\mathcal{L}=\mathcal{O}_S$; assertion (i) then follows immediately from
(\hyperref[I.1.2.2.2]{I, 1.2.2.2}). Likewise, in (ii), we may
restrict to the case where $\mathcal{L}=\mathcal{O}_Z$,
$\mathcal{L}'=\mathcal{O}_{Z'}$, and then the assertion reduces to
Lemma (\hyperref[II.4.3.2.4]{4.3.2.4}).

Let us now prove (i bis). We may assume $Y$ affine
(\hyperref[II.4.6.4]{4.6.4}), hence $\mathcal{L}$ ample
(\hyperref[II.4.6.6]{4.6.6}); as $s$ runs through the union
of the $\Gamma(X,\mathcal{L}^{\otimes n})$ ($n>0$), the $X_s$ form a
basis of the topology of $X$ (\hyperref[II.4.5.2]{4.5.2, a)}), so by
hypothesis the $j^{-1}(X_s)$ form a basis of the topology of $X'$; we
therefore conclude by Lemma
(\hyperref[II.4.6.13.1]{4.6.13.1, (i)}) and by
(\hyperref[II.4.5.2]{4.5.2, a)}) that $j^*(\mathcal{L})$ is ample.

Let us next prove (iii). We may again assume $Y$ and $Y'$ affine
(\hyperref[II.4.6.4]{4.6.4}), from which it follows that the projection
$h:X'\to X$ is affine (\hyperref[II.1.5.5]{1.5.5}). Since
$\mathcal{L}$ is ample (\hyperref[II.4.6.6]{4.6.6}), as $s$
runs through the set of sections of the $\mathcal{L}^{\otimes n}$ ($n>0$)
over $X$ such that $X_s$ is affine, the $X_s$ cover $X$
(\hyperref[II.4.5.2]{4.5.2, a')}); hence the $h^{-1}(X_s)$ are affine
(\hyperref[II.1.2.5]{1.2.5}) and cover $X'$; it therefore follows again
from Lemma (\hyperref[II.4.6.13.1]{4.6.13.1, (i)}) and from
(\hyperref[II.4.5.2]{4.5.2, a')}) that $\mathcal{L}'$ is ample, the
morphism $f_{(Y')}$ being quasi-compact
(\hyperref[I.6.6.4]{I, 6.6.4, (iii)}).

To prove (iv), note first that $f_1\times_S f_2$ is quasi-compact
(\hyperref[I.6.6.4]{I, 6.6.4, (iv)}). We may moreover assume $S$,
$Y_1$ and $Y_2$ affine ((\hyperref[II.4.6.4]{4.6.4}) and
(\hyperref[I.3.2.7]{I, 3.2.7})), hence $\mathcal{L}_i$ ample ($i=1,2$)
(\hyperref[II.4.6.6]{4.6.6}). The open subsets
$(X_1)_{s_1}\times_S(X_2)_{s_2}$ form a covering of
$X_1\times_S X_2$ as $s_i$ runs through the sections of the
$\mathcal{L}_i^{\otimes n_i}$ such that $(X_i)_{s_i}$ is affine
(\hyperref[II.4.5.2]{4.5.2, a')}). Moreover, by replacing $s_1$ and
$s_2$ by suitable powers, which does not change the
$(X_i)_{s_i}$, we can always assume $n_1=n_2$. We then deduce from
(\hyperref[II.4.6.13.1]{4.6.13.1, (ii)}) and from
(\hyperref[II.4.5.2]{4.5.2, a')}) that
$\mathcal{L}_1\otimes_S\mathcal{L}_2$ is ample, whence the assertion, since
$Y_1\times_S Y_2$ is affine (\hyperref[II.4.6.6]{4.6.6}).

It remains to prove (ii). By the same argument as in
(\hyperref[II.4.4.10]{4.4.10}), but using here
(\hyperref[II.4.6.4]{4.6.4}), we may restrict to the case where $Z$ is
affine. Since $\mathcal{K}$ is then ample, and $Y$ quasi-compact, there exist
finitely many sections $s_i\in\Gamma(Y,\mathcal{K}^{\otimes k_i})$
such that the $Y_{s_i}$

\oldpage[II]{93}

are affine and cover $Y$
(\hyperref[II.4.5.2]{4.5.2, a')});
replacing the $s_i$ by suitable powers, we may moreover
assume all the $k_i$ equal to the same integer $k$. Let $s'_i$ be the
sections of $f^*(\mathcal{K}^{\otimes k})$ over $X$ corresponding
canonically to the $s_i$, so that the
$X_{s'_i}=f^{-1}(Y_{s_i})$ (4.6.1.13, (i))
cover $X$. Since $\mathcal{L}|X_{s'_i}$ is ample
((\hyperref[II.4.6.4]{4.6.4}) and
(\hyperref[II.4.6.6]{4.6.6})), there exist for each $i$ finitely many
sections $t_{ij}\in\Gamma(X,\mathcal{L}^{\otimes n_{ij}})$ such that the
$X_{t_{ij}}$ are affine, contained in $X_{s'_i}$, and cover
$X_{s'_i}$ (\hyperref[II.4.5.2]{4.5.2, a')}); we may moreover
assume all the $n_{ij}$ equal to the same number $n$. This being so, $X$ is
separated and quasi-compact, so there exist an integer $m>0$ and, for every
$(i,j)$, a section
\[
  u_{ij}\in\Gamma(X,\mathcal{L}^{\otimes n}\otimes
  f^*(\mathcal{K}^{\otimes mk}))
\]
such that $t_{ij}\otimes s_i'^{\otimes m}$ is the restriction to $X_{s'_i}$
of $u_{ij}$ (\hyperref[I.9.3.1]{I, 9.3.1}); moreover, we have
$X_{u_{ij}}=X_{t_{ij}}$, so the $X_{u_{ij}}$ are affine and cover
$X$. We may moreover assume that $m$ is of the form $nr$; if we set
$n_0=rk$, we therefore see (\hyperref[II.4.5.2]{4.5.2, a')}) that
$\mathcal{L}\otimes_{\mathcal{O}_X}f^*(\mathcal{K}^{\otimes n_0})$ is
ample. Moreover,
there exists $h_0>0$ such that $\mathcal{K}^{\otimes h}$ is generated by its
sections over $Y$ as soon as $h\geq h_0$
(\hyperref[II.4.5.5]{4.5.5}); \emph{a fortiori},
$f^*(\mathcal{K}^{\otimes h})$ is generated by its sections over
$X$ for $h\geq h_0$, by the definition of inverse images
((\hyperref[0.3.7.1]{0, 3.7.1}) and
(\hyperref[II.4.4.1]{4.4.1})). We conclude that
$\mathcal{L}\otimes f^*(\mathcal{K}^{\otimes n_0+h})$ is ample as soon as
$h\geq h_0$ (\hyperref[II.4.5.6]{4.5.6}), which completes the
proof.

\begin{remark}[4.6.14]
\label{II.4.6.14}
Under the conditions of (ii), one should be careful not to believe that
$\mathcal{L}\otimes f^*(\mathcal{K})$ is ample for $g\circ f$; indeed,
since $\mathcal{L}\otimes f^*(\mathcal{K}^{-1})$ is also ample for $f$
(\hyperref[II.4.6.5]{4.6.5}), we would conclude that $\mathcal{L}$ would be
ample for $g\circ f$; taking in particular for $g$ the identity
morphism, every invertible $\mathcal{O}_X$-module would be ample for $f$, which
is not the case in general (see (\hyperref[II.5.1.6]{5.1.6}),
(\hyperref[II.5.3.4]{5.3.4, (i)}) and
(\hyperref[II.5.3.1]{5.3.1})).
\end{remark}

\begin{proposition}[4.6.15]
\label{II.4.6.15}
Let $f:X\to Y$ be a quasi-compact morphism, $\mathcal{J}$ a locally
nilpotent quasi-coherent sheaf of ideals of $\mathcal{O}_X$, $Z$ the
closed sub-prescheme of $X$ defined by $\mathcal{J}$, and $j:Z\to X$
the canonical injection. For an invertible $\mathcal{O}_X$-module
$\mathcal{L}$ to be ample for $f$, it is necessary and sufficient that
$j^*(\mathcal{L})$ be ample for $f\circ j$.
\end{proposition}

Indeed, the question being local on $Y$
(\hyperref[II.4.6.4]{4.6.4}), we may assume $Y$ affine; $X$ being
then quasi-compact, we may assume $\mathcal{J}$ nilpotent. Taking into account
(\hyperref[II.4.6.6]{4.6.6}), the proposition is then none other than
(\hyperref[II.4.5.13]{4.5.13}).

\begin{corollary}[4.6.16]
\label{II.4.6.16}
Let $X$ be a locally Noetherian prescheme and $f:X\to Y$ a quasi-compact
morphism. For an invertible $\mathcal{O}_X$-module $\mathcal{L}$
to be ample for $f$, it is necessary and sufficient that its inverse image
$\mathcal{L}'$ under the canonical injection $X_{\mathrm{red}}\to X$ be ample
for $f_{\mathrm{red}}$.
\end{corollary}

We have already seen that the condition is necessary
(\hyperref[II.4.6.13]{4.6.13, (vi)}); conversely, if it is fulfilled,
we may restrict, in order to prove that $\mathcal{L}$ is ample for $f$, to the
case where $Y$ is affine (\hyperref[II.4.6.4]{4.6.4}); then
$Y_{\mathrm{red}}$ is also affine, so $\mathcal{L}'$ is ample
(\hyperref[II.4.6.6]{4.6.6}), and the same is true of $\mathcal{L}$ by
virtue of (\hyperref[II.4.5.13]{4.5.13}), since then $X$ is
Noetherian and $X_{\mathrm{red}}$ is a closed sub-prescheme of $X$ defined by
a nilpotent quasi-coherent sheaf of ideals
(\hyperref[I.6.1.6]{I, 6.1.6}).

\begin{proposition}[4.6.17]
\label{II.4.6.17}
The notation and hypotheses being those of
(\hyperref[II.4.4.11]{4.4.11}), for $\mathcal{L}''$ to be ample
relative to $f''$, it is necessary and sufficient that $\mathcal{L}$ be ample
relative to $f$ and $\mathcal{L}'$ ample relative to $f'$.
\end{proposition}

\oldpage[II]{94}

The necessity of the condition follows from
(\hyperref[II.4.6.13]{4.6.13, (i bis)}). To see that the condition is
sufficient, we may restrict to the case where $Y$ is affine, and then the fact
that $\mathcal{L}''$ is ample follows from criterion
(\hyperref[II.4.5.2]{4.5.2, a)}) applied to $\mathcal{L}$,
$\mathcal{L}'$ and $\mathcal{L}''$, observing that a section of
$\mathcal{L}$ over $X$ extends (by $0$) to a section of
$\mathcal{L}''$ over $X''$.

\begin{proposition}[4.6.18]
\label{II.4.6.18}
Let $Y$ be a quasi-compact prescheme, $\mathcal{S}$ a quasi-coherent graded
$\mathcal{O}_Y$-algebra of finite type,
$X=\operatorname{Proj}(\mathcal{S})$, and $f:X\to Y$ the structural morphism.
Then $f$ is of finite type, and there exists an integer $d>0$ such that
$\mathcal{O}_X(d)$ is invertible and $f$-ample.
\end{proposition}

By virtue of (\hyperref[II.3.1.10]{3.1.10}), there exists an integer $d>0$
such that $\mathcal{S}^{(d)}$ is generated by $\mathcal{S}_d$. We know that,
under the canonical isomorphism between $X$ and
$X^{(d)}=\operatorname{Proj}(\mathcal{S}^{(d)})$, $\mathcal{O}_X(d)$
is identified with $\mathcal{O}_{X^{(d)}}(1)$
(\hyperref[II.3.2.9]{3.2.9, (ii)}). We therefore see that we are reduced to the case
where $\mathcal{S}$ is generated by $\mathcal{S}_1$; the proposition then follows
from (\hyperref[II.4.4.3]{4.4.3}) and from
(\hyperref[II.4.6.2]{4.6.2}) (taking into account that $f$ is a morphism
of finite type (\hyperref[II.3.4.1]{3.4.1})).

